\documentclass[12pt, a4paper]{scrartcl}
\usepackage{fullpage}
\usepackage[latin1]{inputenc} %Zeichensatzkodierung (auch Sonderzeichen -> �,�,�)
\usepackage[ngerman, english]{babel} %Sprache -> Silbentrennung, automatisch generierte Texte auf deutsch; ngerman -> NEUE deutsche Rechtschreibung
\usepackage[T1]{fontenc} %Sonderzeichen allg
\usepackage{lmodern} %PDF-optimierte Schrift

\usepackage{amsmath, amssymb, amsbsy, amsthm} %Mathepakete ; amsthm: Theorem-Umgebung; muss nach amsmath eingebunden werden
\usepackage{eqnarray}
\usepackage[x11names]{xcolor} % muss vor "`mcode"' da dieses bereits intern das color-Paket l�dt und es dann einmal mit und einmal ohne Option geladen wird, was nicht geht 
\usepackage{graphicx} % Grafiken einf�gen
%\usepackage{listings} % zum Einbinden von Code, ist bereits in "`mcode"' enthalten
\usepackage{enumerate}
\usepackage{tabularx} % erwiterte Tabelleneigenschaften
\usepackage{here} % erzwingen der Position von Float-Objekten
\usepackage{bibgerm} % F�r das Literaturverzeichnis
% \usepackage{extarrows} % Pfeile mit Beschriftung, deren L�nge automatisch angepasst wird
\usepackage{algorithm2e} % Pseudocode
\usepackage{algpseudocode} % Pseudocode
\usepackage{bm} % fett schreiben in Matheumgebung

\usepackage[onehalfspacing]{setspace}
% um Zeilenabstand zu setzen

%\usepackage{bibgerm} % F�r das Literaturverzeichnis



%Definition der gew�nschten Theorem-Stile:
\newtheorem{theorem}{Theorem}[section]
\newtheorem{lemma}[theorem]{Lemma} %hier [section] nicht noch einmal wiederholen, da durch [satz] die Nummerierung schon genauo wie in "Satz" definiert
\newtheorem{corollary}[theorem]{Corrolary}

\theoremstyle{definition}
\newtheorem{definition}[theorem]{Definition}
\newtheorem{convention}[theorem]{Convention}
\newtheorem{example}[theorem]{Example}

\newtheoremstyle{my_remark}
{}{}{}{}{\itshape}{:}{.5em}{}

\theoremstyle{my_remark}
\newtheorem*{remark}{Remark}

%\newenvironment{proof}{\begin{proof}[Proof:]}{\end{proof}}


\newcommand{\R}{\mathbb{R}}
\newcommand{\N}{\mathbb{N}}
\newcommand{\blangle}{\big\langle}
\newcommand{\brangle}{\big\rangle}
\DeclareMathOperator*{\argmin}{arg\,min}
\DeclareMathOperator*{\argmax}{arg\,max}
\newcommand{\erl}{\textcolor{Green3}{ - erledigt}}

\allowdisplaybreaks[1] % [1]: Umgebungen wie align d�rfen Seitenumbr�che enthalten

%\setlength{\voffset}{-28.4mm}
%\setlength{\hoffset}{-1in}
%\setlength{\topmargin}{20mm}
%\setlength{\oddsidemargin}{25mm}
%\setlength{\evensidemargin}{25mm}
%\setlength{\textwidth}{160mm}
%
\setlength{\parindent}{0pt}
\setlength{\parskip}{1ex}
%
%\setlength{\textheight}{235mm}
%\setlength{\footskip}{20mm}
%\setlength{\headsep}{50pt}
%\setlength{\headheight}{0pt}



\begin{document}

%TODO
 %% check command for ,...,

%% Titelseite
\thispagestyle{empty}       % keine Seitennummer
\parbox{1.5cm}{\includegraphics[height=1.5cm]{Pictures/TUMLogo_schwarz}}\hfill
\parbox{1.5cm}{\includegraphics[height=1.5cm]{Pictures/Mathe_schwarz}}%
\vspace*{1.5cm}
\begin{center}

{\Large Technische Universit�t M�nchen} \\
\vspace*{1cm}
{\large \textsc{Department of Mathematics}} 

\vspace*{3cm}
\begin{doublespace}
{\Huge \textbf{Bundle Methods for Nonsmooth, Nonconvex Problems with Inexact Information}} \\
\vspace*{2.5cm}
{\large Master's Thesis}\linebreak 
{\large by}\linebreak
{\large Annika Stegie}
\end{doublespace}
\end{center}

\vfill
 
\begin{tabular}{ll}
Supervisor: & Prof. Dr. Michael Ulbrich\\[1.5ex]
Advisers:   & Dr. Andre Milzarek \\
            & Lukas Hertlein (M.Sc.) \\[1.5ex]
Submission date: & August 31, 2017
\end{tabular}
\newpage    % Seitenwechsel
%\textcolor{red}{second title page} \\
%\textcolor{red}{Author and title on back of thesis}
%%Selbsst�ndigkeitserkl�rung
%\thispagestyle{empty}

\vspace*{150mm}
I hereby declare that this thesis is my own work and that no other sources have been used except those clearly indicated and referenced.
\vspace{2cm}

Garching, \textcolor{red}{Date}

%\section*{Acknowledgments}

First of all I thank Prof. Ulbrich for the opportunity to write my thesis at his chair.
I thank my advisers for the the interesting topic and the good support and guidance in the course of this thesis. Particularly I want to thank Dr. Andre Milzarek for the good communication even over longer distances and Lukas Hertlein for taking over the mentoring so naturally.


\thispagestyle{empty}

%\section*{List of Symbols}

\thispagestyle{empty}

\begin{tabular}{c@{}c@{} l@{}}
	\(x^k\) && iterates \\
	\(\hat{x}^k\) && current stability center \\
	\(f(x)\) && exact evaluation of \(f\) at point \(x\) \\
	\(f_k\) && value of \(f\) at point \(x^k\) from the oracle, may be inexact \\
	\(\hat{f}_k\) && value of \(f\) at stability center \(\hat{x}^k\) from the oracle, may be inexact \\
	\(g^k\) && a subgradient of \(f\) at \(x^k\) (can be exact or inexact) \\
	\(s^k\) && a subgradient of the convexified objective at \(x^k\) (can be exact or inexact) \\
	\(m_k\) && cutting plane model of the objective function \(f\) \\
	\(M_k\) && cutting plane model of the convexified objective function \\
	\(e_j^k\) && linearization error \\
	\(c_j^k\) && linearization error of the convexified objective function \\
	\(\eta_k\) && convexification parameter \\
	\(\gamma\) & \(>0\) & safeguarding parameter for \(\eta\)-calculation \\
  \(\alpha_j^k\) && Lagrange multipliers of the subproblem \\
	\(d^k\) && minimizer of the subproblem \\
	\(\xi\) && variable from reformulation of the subproblem \\
	\(J_k\) && index set at iteration \(k\) \\
	\(t_k\) &\(>0\)& prox-parameter \\
	\(\delta_k\) && nominal decrease \\
	\(m\) &\(\in (0,1) \quad\)& decrease parameter \\
	\(G^k\) && aggregate subgradient \\
	\(S^k\) && aggregate subgradient of convexified objective \\
	\(E^k\) && aggregate error \\
	\(C^k\)	&& aggregate error of convexified objective \\
	\(\rho\) & \(\in (0,\frac{1}{2}) \quad \) & correction parameter \\
	K & \(>0\) & control parameter for length of direction vector \\ 
	\(\theta_k\) && length control factor
\end{tabular}


\setcounter{tocdepth}{3} % Im Inhaltsverzeichnis werden nur section und subsection ber�cksichtigt, keine tieferen Ebenen
\thispagestyle{empty} % keine Seitenzahl
\addtocontents{toc}{%
  \protect\contentsline{section}{List of Symbols}{}}
\tableofcontents
\newpage

\pagenumbering{arabic}
%\section{Introduction}

There exists a sound and board theory of classical nonlinear optimization. However, this theory puts strong differentiability requirements on the given problem. Requirements that cannot always be fulfilled in practice.
Examples for such practical application reach from problems in physics and mechanical engineering \cite{Clarke1990} over optimal control problems up to data analysis \cite{Bagirov2014} and machine learning \cite{Smola2007}.
Other possible fields of applications are risk management and financial calculations \cite{Nesterov2016,Teo2010}. 
Additionally there exist so called stiff problems that are theoretically smooth but numerically nonsmooth due to rapid changes in the gradient \cite{Maekelae1992}. \\
There exists therefore a need for nonsmooth, that is not necessarily differentiable, optimization algorithms. A lot of the underlying theory and was developed in the 1970's, also driven by the ``First World Conference on Nonsmooth Optimization'' taking place in 1977 \cite{Mifflin2012}.
Now, there exists a well understood theoretical framework of nonsmooth analysis to create the basis for practical algorithms. \\
The most popular methods to tackle nonsmooth problems at the moment are bundle methods \cite{Hare2016}. First developed only for convex functions \cite{Lemarechal1978} the method was soon extended to cope also with nonconvex objective functions \cite{Mifflin1982}. \\
Some time later these algorithms were again enhanced to deal with inexact information of the function value, the subgradient or both.\\
Some natural applications for these cases are derivative free optimization and stochastic simulations \cite{Hare2016}. \textcolor{red}{Some more examples from different sources? Bilevel Problems?} \\


The basic idea of bundle methods is to model the original problem by a simpler function, often some sort of stabilized cutting plane model, that is minimized as a subproblem of the algorithm \cite{Hiriart-Urruty1993}. 

\textcolor{orange}{Adapt this part to what I finally really do: \\
In this thesis two different types of model functions will be examined that allow the use of inexact information in small to medium-scale problems as well as in large-scale problems. A limited memory approach is examined for the latter case. \\}
\textcolor{red}{what new? Combination of large-scale and inexact information - why needed \\
don't forget What - why - how}

\textcolor{orange}{Adapt this part to what I finally really do: \\
This thesis is organized as follows: \\
introduction of the most important definitions and results of nonsmooth analysis. Then the introduction of a very basic bundle algorithm which is then generalized for nonconvex functions with nonsmooth optimization. \\
Throughout study of this algorithm including comparison to other approaches to tackle inexact information. \\
Introduction of variable metric (bundle) algorithm to tackle large-scale applications. ``discussion'' how far this is compatible with inexactness. \\
Numerical testing \\ 
discussion}

This thesis is written with the acedemic 'we'.


First a proximal bundle method \textcolor{red}{Difference between different regularizations explained before}...
large-scale optimization: a metric bundle method instead of a proximal bundle method -> limited memory approach \\


\textcolor{red}{from PhD-thesis}
\begin{itemize}
	\item 
\end{itemize}

%\section{Preliminaries}
\label{sec_prelim}

When it comes to nonsmooth objective functions the  derivative based framework of nonlinear optimization methods does not work any more. Meanwhile there exists a well understood theory of 'subdifferential calculus' that gives similar results in the nondifferentiable case.
The most important definitions and results of this theory together with some remarks on notation are stated in this chapter.


\subsection{Notation}

Let \(x\) denote a column vector. The transpose of \(x\) is denoted by \(x^{\top}\). The scalar product is written \(\langle \cdot, \cdot \rangle\).
In this thesis generally the euclidean norm is used and denoted by \(\|\cdot\|\). In chapter \ref{sec_variable_metric} additionally a norm is used that is induced by a symmetric matrix. Here we use the notation \(\|x\|^2_A = {\langle x, Ax \rangle}\).
Inequalities written for vectors \(x^1 \leq x^2\), \(x^1, x^2 \in \R^n\) are to be read component wise.
With \(0\) we denote the zero vector of appropriate size.  The identity matrix of appropriate size is written as \(\mathbb{I}\).

As we work with numerical methods in this thesis occur a lot of sequences of various dimensions. For vectors iteration indices are indicated by a superscript \(x^k\) whereas the components are  indicated by subscripts \(x = \left(x_1,x_2,...,x_n\right)^{\top}\). Sequences of numbers and matrices a indexed with subscripts.
For (sub-)sequences where \(k\) comes from an index set \(K \subset \N\) we write \(\{x^k\}_{k \in K}\). If \(k\) is in the natural numbers this notation is shortened to \(\{x^k\}\).
We denote the open ball around \(x\) with radius \(r\) with \(B_{r}(x)\).
The subset relation is denoted by \(A\subset B\). It is to be read in the sense that \(A\) is a subset of \(B\) or that \(A=B\). 

%\textcolor{red}{
%\begin{itemize}
	%\item iteration index as superscript \(x^k\), entry index as subscript \(x_i\)
	%\item  is the scalar product
	%\item more?
%\end{itemize}}


%\textcolor{red}{Theoretical Background, nonsmooth Analysis ???}

%\textcolor{red}{Check if requirements on functions are stated and defined.}\\

\subsection{Nonsmooth Analysis}

Throughout this thesis we consider different optimization problems of the form

\begin{equation*}
\min_{x} f(x) \quad \text{s.t.} \quad x \in X \subset \R^n
\label{min_prob}
\end{equation*}

where \(f\) is a possibly nonsmooth function.

%Therefore the most important definitions and results needed when working with nonsmooth functions are stated in this section.

Nonsmooth functions have kinks where a unique gradient cannot be defined. It is however possible to define a set of tangents to the graph called subdifferential.
The subdifferential was first defined for convex functions.

\begin{definition}[{\cite[Definition 1.2.1, p. 241]{Hiriart-Urruty1996}}]
\label{conv_subdiff}
	Let \(f: \R^n \to \R \) be a convex function. The \emph{subdifferential} of \(f\) at \(x\in \R^n\) is the set 
	\[ \partial f(x) := \left\{g \in \R^n \mid f(y)-f(x) \geq \Langle g,y-x \Rangle \quad \forall y \in \R^n \right\}. \]
\end{definition}

The subdifferential is a set valued mapping. It is closed and convex. If \(f\) is differentiable, its subdifferential is single valued and coincides with its gradient \(\partial f(x) = {\nabla f(x)}\) \cite{Rockafellar1970}.


It is also possible to define a subdifferential for nonconvex functions. This is the subdifferential we will work with in this thesis most of the time.

\begin{definition}[c.f. {\cite[p. 25, 27]{Clarke1990}}]
\label{nonconv_subdiff}
	Let \(f: \R^n \to \R\) be locally Lipschitz (and not necessarily convex).
	The \emph{subdifferential} of \(f\) at \(x \in \R^n\) is the set
\[ \partial f(x) := \left\{g \in \R^n \biggm| \limsup_{y \to x,~ h \searrow 0}\frac{f(y+hv)-f(y)}{h} \geq \Langle g, v\Rangle \quad \forall v \in \R^n \right\}. \]
\end{definition} 

All convex functions are locally Lipschitz \cite[Theorem 3.1.1, p. 16]{Hiriart-Urruty1996} so the above definition holds also for convex functions. In fact if the function is convex the subdifferential from definition \ref{nonconv_subdiff} is equivalent to the one from definition \ref{conv_subdiff} \cite[Proposition 2.2.7, p. 36]{Clarke1990}.
Due to this equivalence we call elements from both subdifferentials subgradients.

\begin{remark}
	It is important to observe that subgradient inequality 
	\begin{equation}
	f(y)-f(x) \geq \Langle g,y-x \Rangle \quad \forall y \in \R^n 
\label{subgr_inequ}
\end{equation}
	only holds in the convex case.
\end{remark}

There is also a sum rule for the subdifferential.

\begin{proposition}[{\cite[Proposition 2.3.3, p. 38]{Clarke1990}}]
	Let \(F(x)=\sum_i{f_i(x)}\) be a finite sum of nondifferentiable functions. Then it holds
	
	\[ \partial F(x) \subset \sum_i{\partial f_i(x)}. \]
\end{proposition}

Analogous to the \(\mathcal{C}^1\)-case some first order optimality conditions can be stated.
For nondifferentiable functions a \emph{stationary point} \(x\) of the function \(f\) is characterized by \cite[p. 38]{Clarke1990}
\[ 0 \in \partial f(x). \]
If the function \(f\) is convex, then every stationary point is a minimum.

%\textcolor{red}{-> characterization of minimum and stationary point}

A drawback of the subdifferential is that it does not indicate how near the evaluated point is to a stationary point or minimum of a function. This can only be seen if the evaluated point is already stationary.

This issue is addressed by the \(\varepsilon\)-\emph{subdifferential}. It gathers all information in a small neighborhood of the point \(x\).

For convex functions an \(\varepsilon\)-\emph{subgradient} of \(f(x)\) is defined as a vector \(g \in \R^n\) satisfying the inequality

\[ f(y)-f(x) \geq \Langle g,y-x \Rangle - \varepsilon \quad \forall y \in \R^n.\]

The \(\varepsilon\)-subdifferential is then the set
\[ \partial_{\varepsilon} f(x) := \left\{g \in \R^n \mid g \text{ is an } \varepsilon \text{-subgradient of } f(x) \right\}. \]

For nonconvex functions the subdifferential that is used in this thesis is the \emph{Fr\'{e}chet} \(\varepsilon\)-\emph{subdifferential}.

\begin{definition}[c.f. {\cite[p. 73]{Jofre1998}}]
	The Fr\'{e}chet \(\varepsilon\)-subdifferential of \(f(x)\) is
	\[ \partial_{[\varepsilon]}f(x) := \left\{ g \in \R^n \biggm| \liminf_{\|h\| \to 0} \frac{f(x+h)-f(x)-\langle g,h\rangle}{\|h\|}  \geq -\varepsilon\right\}. \]
\end{definition}

For \(\varepsilon = 0\) this is called \emph{Fr\'echet subdifferential}.
For convex functions the Fr\'echet \(\varepsilon\)-subdifferential and the \(\varepsilon\)-subdifferential are \emph{not} the same.

In the course of this thesis we sometimes derive stronger results for a smaller class of nonsmooth functions. Those functions are called lower-\(\mathcal{C}^2\) functions and can be defined as follows.

\begin{definition}[c.f. Definition 10.29, p. 447 in {\cite{Rockafellar2009}}]
\label{def_low_C2}
	A function \(f:\R^n \to \R\) is \emph{lower-}\(\mathcal{C}^2\), if on some neighborhood \(\Omega\) of each \(\bar{x} \in \R^n\) there exists a representation 
	
	\[ f(x) = \max_{t\in T}f_t(x) \]
	
	where the functions \(f_t\) are of class \(\mathcal{C}^2\) on \(\Omega\) and \(f_t(x)\) and all its first and second partial derivatives depend continuously on both \(x\) and \((x,t) \in \Omega \times T\).
	The index set \(T\) is a compact space.
\end{definition}

Lower-\(\mathcal{C}^2\) functions are locally Lipschitz continuous \cite[Theorem 10.31, p. 448]{Rockafellar2009}.
%\subsection{Nonsmooth analysis and optimization}

%A necessary assumption on the objective function \(f\) is that it is locally Lipschitz. This assumption assures the well-definedness of the following generalizations of derivatives. \\
%\begin{definition} \cite{Mifflin1977}
	%A function \(f: \R^n \to \R\) is called \emph{locally Lipschitz}  if it is Lipschitz on each bounded subset \(B \subseteq \R^n\)
	%\[ |f(y) - f(x)| \leq C \|y-x\| \quad \forall x,y \in B, \quad C >0. \]
%\end{definition}
%All convex functions are locally Lipschitz \cite{Hiriart-Urruty1996}. \\

%For convex functions one can define so called subgradients as a generalization of the usual derivative. They are defined using the  directional derivative. \\
%\begin{definition} \cite{Hiriart-Urruty1996}
	%The \emph{directional derivative}  of a convex function \(f\) at \(x\) in direction \(d\) is 
	%\[ f'(x,d) := \lim_{\lambda \downarrow 0}\frac{f(x+\lambda d) - f(x)}{\lambda}. \]
%\end{definition}

%\begin{definition} \cite{Hiriart-Urruty1996}
	%Let \(f\) convex. The \emph{subdifferential} \(\partial f(x)\) of \(f\) at \(x\) is the nonempty compact convex set
	%\[ \partial f(x) = \{g \in \R^n | f'(x,d) \geq \langle g,d\rangle \forall d \in \R^n \}. \]
%\end{definition}
%The subdifferential is a convex set, that supports the graph of the function \(f\) from below. If \(f\) is differentiable at the point \(x\), the subdifferential reduces to the gradient at that point \cite{Hiriart-Urruty1996}. \\

%This concept was generalized by Clarke for nonconvex functions. First a generalization of the directional derivative is given: 
%\begin{definition} \cite{Clarke1990}
	%Let \(f: \R^n \to \R \) locally Lipschitz. The \emph{generalized directional derivative}  of \(f\) at \(x\) in direction \(d\) is given by
	%\[ f^{\circ}(x,d):= \limsup_{\substack{y \to x \\ \lambda \downarrow 0}} \frac{f(y+\lambda d)-f(y)}{\lambda}. \]
%\end{definition}

%This allows for the following definition.
%\begin{definition} \cite{Clarke1990}
	%The \emph{generalized gradient} of the locally Lipschitz function \(f\) at \(x\) is a nonempty convex compact set \(\partial f(x)\) given by
	%\[ \partial f(x) := \{g \in \R^n | f^{\circ}(x,d) \geq \langle g,d\rangle \forall d \in \R^n \}. \]
%\end{definition}

%\textcolor{red}{!!!other definition --> take definition from rockefellar/Hare directly from Paper!!!}

%If \(f\) is a convex function the generalized gradient coincides with the subdifferential \(\partial f\) of \(f\) \cite{Clarke1990}.


%\textcolor{red}{For the biggest part of this thesis and if not otherwise stated \(X = \R^n\).} \\
%\textcolor{red}{Also write something about inexactness? specify \(X\) more precisely? Convex?} \\
%\textcolor{red}{Just definition, lemma, theorem or a bit explanation around it?\\
%better just in Text without Definition, ...} \\
%\textbf{\textcolor{red}{See if requirements in definitions and theorems meet what is needed/provided later.}} \\

%\begin{definition} % as in 
		%The function \(f: \R^n \to \R\) is called \emph{Lipschitz near \(x \in \R^n\)}\cite{Clarke1990} if there exist \(C > 0\) and \(\varepsilon > 0\) such that 
	%\[ |f(y_2) - f(y_1)| \leq C\|y_2 -y_1\| \quad \forall y_1, y_2 \in \mathbf{B}_{\varepsilon}(x). \]
%\end{definition}


%Throughout this thesis the objective function is assumed to be locally Lipschitz. 

%For convex functions one can define so called subgradients as a generalization of the usual derivative. They are defined using the  directional derivative. \\
%\begin{definition} \cite{Hiriart-Urruty1996}
%	The \emph{directional derivative}  of a convex function \(f\) at \(x\) in direction \(d\) is 
%	\[ f'(x,d) := \lim_{\lambda \downarrow 0}\frac{f(x+\lambda d) - f(x)}{\lambda}. \]
%\end{definition}

%\begin{definition} \cite{Hiriart-Urruty1996}
	%Let \(f\) convex. The \emph{subdifferential} \(\partial f(x)\) of \(f\) at \(x\) is the nonempty compact convex set
	%\[ \partial f(x) = \{g \in \R^n | f'(x,d) \geq \langle g,d\rangle \forall d \in \R^n \}. \]
%\end{definition}
%The subdifferential is a convex set, that supports the graph of the function \(f\) from below. If \(f\) is differentiable at the point \(x\), the subdifferential reduces to the gradient at that point \cite{Hiriart-Urruty1996}. \\
%
%This concept was generalized by Clarke for nonconvex functions. First a generalization of the directional derivative is given: 
%\begin{definition} \cite{Clarke1990}
	%Let \(f: \R^n \to \R \) locally Lipschitz. The \emph{generalized directional derivative}  of \(f\) at \(x\) in direction \(d\) is given by
	%\[ f^{\circ}(x,d):= \limsup_{\substack{y \to x \\ \lambda \downarrow 0}} \frac{f(y+\lambda d)-f(y)}{\lambda}. \]
%\end{definition}
%
%This allows for the following definition.
%\begin{definition} \cite{Clarke1990}
	%The \emph{generalized gradient} of the locally Lipschitz function \(f\) at \(x\) is a nonempty convex compact set \(\partial f(x)\) given by
	%\[ \partial f(x) := \{g \in \R^n | f^{\circ}(x,d) \geq \langle g,d\rangle \forall d \in \R^n \}. \]
%\end{definition}
%
%\textcolor{red}{!!!other definition --> take definition from rockefellar/Hare directly from Paper!!!}

%If \(f\) is a convex function the generalized gradient coincides with the subdifferential \(\partial f\) of \(f\) \cite{Clarke1990}.

%\textcolor{red}{Why epsilon Subdifferential?\\
%implementable stopping criterion \\
%dual form of bundle algorithms = same as stopping criterion?}

%\begin{definition} \cite{Mifflin1977}
	%The function \(f: \R^n \to \R\)  is \emph{semismooth} at \(x \in \R^n\)  if \(f\) is Lipschitz on a ball \(\mathbb{B}_\varepsilon(x)\)around x and for each \(d \in \R^n\) and for any sequences \(\{t_k\} \subseteq \R_{+}\), \(\{\theta_k\} \subseteq \R^n\) and \(\{g_k\} \subseteq \R^n\) such that
	%\[ \{t_k\} \downarrow 0, \quad \{\theta_k/t_k\} \to 0 \in \R^n \quad \text{and} \quad g_k \in \partial f(x+t_kd+\theta_), \]
	%the sequence  \(\{\langle g_k,d\rangle\}\) has exactly one accumulation point.
%\end{definition}

%\textcolor{red}{check if I need \(\infty\)-functions and if something changes then}
%\textcolor{red}{\(\epsilon\)-subdifferential \\
%continuity properties of generalized gradient? \\
%definitions from chapter inexact information \\
%definition of inexactness for nonconvex kind of generalization of \(\epsilon\)-subdifferential for nonconvex case (Noll, inex, nonconv)}

%\section{A Basic Bundle Method}
\label{sec_conv_ex}

When bundle methods were first introduced in 1975 by Claude Lemar�chal and Philip Wolfe they were developed to minimize a convex (possibly nonsmooth) function \(f\) for which at least one subgradient at any point \(x\) can be computed \cite{Mifflin2012}.
To provide an easier understanding of the proximal bundle method in \cite{Hare2016} and stress the most important ideas of how to deal with nonconvexity and inexactness first a basic bundle method is shown here. 

%\textcolor{red}{link to chapter?}

Bundle methods can be interpreted in two different ways: From the dual point of view one tries to approximate the \(\varepsilon\)-subdifferential to finally ensure first order optimality conditions. The primal point of view interprets the bundle method as a stabilized form of the cutting plane method where the objective function is modeled by tangent hyperplanes \cite{Hare2010}. We focus here on the primal approach.

%\textcolor{blue}{notation, definitions}\\

%\textcolor{red}{already done in previous preliminaries chapter?}

%\subsection{A basic bundle method}

%\textcolor{red}{In the next two sections the function \(f\) is assumed to be convex.} \\

\subsection{Derivation of the Bundle Method}
This section gives a short summery of the derivations and results of chapter XV in \cite{Hiriart-Urruty1993} where a primal bundle method is derived as a stabilized version of the cutting plane method. If not otherwise indicated the results in this section are therefore taken from \cite{Hiriart-Urruty1993}.

The optimization problem considered in this section is 
\begin{equation}
  \min_{x} f(x) \quad \text{s.t.} \quad x \in X
\label{min_prob_basic}
\end{equation}
where \(f\) is a convex but possibly nondifferentiable function and \(X \subseteq \R^n\) is a  closed and convex set.

%\textcolor{red}{Define Problem again?? Incorporate ``set-constraint'' by writing \(h(x):= f(x)+\mathtt{i}_X\). \(\rightarrow\) later???}

%\textcolor{blue}{explanation}

\subsubsection{A Stabilized Cutting Plane Method}

The geometric idea of the \emph{cutting plane method} is to build a piecewise linear model of the objective function \(f\) that can be minimized more easily than the original objective function.
This model is built from a \emph{bundle} of information that is gathered in the previous iterations.
In the \(k\)'th iteration, the bundle consists of the previous iterates \(x^j\), the respective function values \(f(x^j)\) and a subgradient at each point \(g^j \in \partial f(x^j)\) for all indices \(j\) in the index set \(J_k\). From each of these triples, one can construct a linear function 

\begin{equation*}
	l_j(x) = f(x^j) + \Langle g^j,x-x^j \Rangle
\label{lin_fun}
\end{equation*}  

where \(f(x^j) = l_j(x^j)\) and due to convexity \(f(x) \geq l_j(x), ~ x \in X\).\\ % x \in \R^n??
The objective function \(f\) can then be approximated by the piecewise linear function

\begin{equation*}
	m_k(x) = \max_{j \in J_k} l_j(x).
\label{cp_model}
\end{equation*}

A new iterate \(x^{k+1}\) is found by solving the subproblem

\begin{equation*}
	\min_{x} m_k(x) \quad \text{s.t.} \quad x \in X.
\label{cp_model_fun}
\end{equation*}

\textcolor{red}{Picture of function and cutting plane approximation of it}

This subproblem should of course be easier to solve than the original task. A question that depends a lot on the structure of \(X\). If \(X = \R^n\) or a polyhedron, the problem can be solved easily. Still there are some major drawbacks to the idea. For example if \(X = \R^n\) the solution of the subproblem in the first iteration is always \(-\infty\).
In general we can say that the subproblem does not necessarily have a solution.
To tackle this problem a penalty term is introduced to the subproblem. It then reads

\begin{equation}
	\min \tilde{m}_k(x) =  m_k(x) + \frac{1}{2 t_k}\|x-x^k\|^2 \quad \text{s.t.} \quad x \in X, ~t_k > 0.
\label{stabil_model_fun}
\end{equation}

This new subproblem is strongly convex and therefore always has  a unique solution.
%\textcolor{red}{how much explanation here?}
% \(\max_{j \in J_k} l_j(\hat{x}^k + d)\)

%\textcolor{red}{Some nice sentences to explain the term a little bit more and to lead over to the next paragraph.\\ 
%To understand the deeper motivation of this term see \cite{Hiriart-Urruty1993}. For this introduction it suffices to see that due to the regularization term the subproblem is now strongly convex and therefore always uniquely solvable.} \\

The regularization term can be motivated and interpreted in many different ways, c.f. \cite{Hiriart-Urruty1993}. From different possible regularization terms the most popular in bundle methods is the penalty-like regularization used here.

The second major step towards the bundle algorithm is the introduction of a so called \emph{stability center} or  \emph{serious point} \(\hat{x}^k\). It is the iterate that yields the ``best'' approximation of the optimal point up to the \(k\)'th iteration (not necessarily the best function value though).
The updating technique for \(\hat{x}^k\) is crucial for the convergence of the method: If the next iterate yields a decrease of \(f\) that is ``big enough'', namely bigger than a fraction of the decrease suggested by the model function for this iterate, the stability center is moved to that iterate. If this is not the case, the stability center remains unchanged.

In practice this looks the following:
Define first the \emph{model decrease} \(\delta^M_k\) which is the decrease of the model for the new iterate \(x^{k+1}\) compared to the function value at the current stability center \(\hat{x}^k\)

\begin{equation}
	\delta^M_{k} = f(\hat{x}^k) - m_k(x^{k+1}) \geq 0.
\label{mod_dec}
\end{equation}

%\textcolor{red}{The nominal decrease is in fact stated a little differently for different versions of the bundle algorithm, this is why I added the constant \(a_k \in \R\) here for generalization. In practice the difference between the decreases is not influencing the algorithm as \(\delta_k\) is weighted by the constant \(m \in (0,1)\) for the descent test which compensates \(a_k\).} 

If the actual decrease of the objective function is bigger than a fraction of the model decrease 

\[ f(\hat{x}^k) - f(x^{k+1}) \geq m \delta^M_k, \quad m \in (0,1) \]

set the stability center to \( \hat{x}^{k+1} = x^{k+1}\). This is called a \emph{serious} or \emph{descent step}.
If this is not the case a \emph{null step} is executed and the serious iterate \(\hat{x}^{k+1} = \hat{x}^k\) remains the same .

Besides the model decrease other forms of decrease measures and variations of these are possible. Some are presented in \cite{Hiriart-Urruty1993} and \cite{Oliveira2014}.

\subsubsection{Subproblem Formulations}
The subproblem to be solved to find the next iterate can be rewritten as a smooth optimization problem. For convenience we first rewrite the affine functions \(l_j\) with respect to the stability center \(\hat{x}^k\).

%\textcolor{red}{citation for this???!!!}

\begin{align*}
	l_j(x) &= f(x^j) +  \Langle g^j,x - x^j\Rangle \\
	&= f(\hat{x}^k)+\Langle g^j,x-\hat{x}^k\Rangle-\left(f(\hat{x}^k)-f(x^j) + \Langle g^j,x^j-\hat{x}^k\Rangle\right) \\
	&= f(\hat{x}^k) + \Langle g^j,x - \hat{x}^k\Rangle - e^k_j
\end{align*}	

where
\begin{equation*}
	e^k_j := f(\hat{x}^k)-f(x^j) + \Langle g^j,x^j-\hat{x}^k\Rangle \geq 0 \quad \forall j \in J_k
\label{lin_err}
\end{equation*}

is the \emph{linearization error}. Its nonnegativity property is essential for the convergence theory and will also be of interest when moving on to the case of nonconvex and inexact objective functions. \\


Subproblem (\ref{stabil_model_fun}) can now be written as

\begin{align}
	& \min_{\hat{x}^k+d \in X} \tilde{m}_k(\hat{x}^k+d) = f(\hat{x}^k) + \max_{j \in J_k} \left\{\Langle g^j,d\Rangle - e^k_j\right\} + \frac{1}{2t_k}\|d\|^2
	\label{sub_prob_long}\\
	\Leftrightarrow \quad &\min_{\substack{\hat{x}^k+d \in X, \\ \xi \in \R}} \xi + \frac{1}{2t_k}\|d\|^2 \quad \text{s.t.} \quad \Langle g^j, d\Rangle - e^k_j - \xi \leq 0, \quad j \in J_k
\label{sub_prob_short}
\end{align}


where \(d:= x-\hat{x}^k\) and the constant term \(f(\hat{x}^k)\) was discarded for the sake of simplicity. \\
If \(X\) is a polyhedron this is a quadratic optimization problem that can be solved using standard methods of nonlinear optimization. The pair \((\xi_k, d^k)\) solves (\ref{sub_prob_short}) if and only if

\begin{align}
	d^k &\text{ solves the original subproblem (\ref{sub_prob_long})  and } \nonumber \\
	\xi_k & =\max_{j \in J_k}{g^j}^{\top}d^k - e_j^k = m_k(\hat{x}^k+d^k) - f(\hat{x}^k). \label{xi}
\end{align}
	
The new iterate is given by \(x^{k+1} = \hat{x}^k + d^k\). %\textcolor{red}{citation!!!}\\

\subsection{The Prox-Operator}
\label{subsec_prox_op}

The constraint \(\hat{x}^k+d \in X\) can also be incorporated directly in the objective function by using the indicator function

\[ \mathtt{i}_X(x) = \left\{ \begin{array}{cl} 0, & \text{if } x \in X \\ +\infty, & \text{if } x \notin X \end{array} \right. .\]

This function is convex if and only if the set \(X\) is convex \cite{Rockafellar1996}.

Subproblem (\ref{stabil_model_fun}) then reads 
\begin{equation*}
	\min_{x \in \R^n} m_k(x) + \mathtt{i}_X(x) + \frac{1}{2t_k}\|x-\hat{x}^k\|^2 % \quad \text{s.t.} \quad {g^j}^{\top}d - e^k_j - \xi \leq 0, \quad j \in J_k
\label{sub_prob_fin}
\end{equation*}

with respect to the serious point \(\hat{x}^k\).
%\textcolor{red}{check if f also not put into subproblem before}

The subproblem is now written as the \emph{Moreau-Yosida regularization} of \(\check{f}:= m_k(x) + \mathtt{i}_X(x)\). The emerging mapping is also known as \emph{proximal point mapping} \cite{Hare2010} or \emph{prox-operator}

	\begin{equation}
		prox_{t,\check{f}}(x) = \argmin_{y \in \R^n}\left\{\check{f}(y) + \frac{1}{2t}\|x-y\|^2\right\}, \quad t > 0.
	\label{prox_op}
	\end{equation}

This special form of the subproblems gives the primal bundle method its name, \emph{proximal bundle method}. The mapping also plays a key role when the method is generalized to nonconvex objective functions and inexact information.


\subsection{Aggregate Objects}

We look again at a slightly different formulation of the bundle subproblem 

\begin{align*}
	\min_{\substack{d \in \R^n, \\ \xi \in \R}} & \quad \xi + \mathtt{i}_X+\frac{1}{2t_k}\|d\|^2 \\
	\text{s.t.} & \quad \Langle g^j,d\Rangle - e^k_j - \xi \leq 0, \quad j \in J_k. 
%\label{sub_prob_smooth_iX}
\end{align*}

As the objective function is still convex (\(X\) is a convex set) the following Karush-Kuhn-Tucker (KKT) conditions have to be valid for the minimizer \(\left(\xi_k,d^k\right)\) of the above subproblem \cite{Hiriart-Urruty1996} assuming a constraint qualification holds if the constraint set \(X\) makes it necessary \cite{Solodov2011}.

There exist a subgradient \(\nu^k \in \partial \mathtt{i}_{X}(x^{k+1})\) and Lagrangian multipliers \(\alpha_j, ~j \in J^k\) such that

\begin{align}
	& 0 = \nu^k + \frac{1}{t_k}d^k + \sum_{j \in J^k}\alpha_j g^j \label{KKT_1}\\
	& \sum_{j \in J_k}\alpha_j = 1, \label{KKT_2}\\
	& \alpha_j \geq 0, ~j \in J^k, \label{KKT_3}\\
	& \Langle g^j,d^k\Rangle - e^k_j - \xi_k \leq 0, \label{KKT_4}\\
	& \sum_{j \in J^k} \alpha_j \left(\Langle g^j,d^k\Rangle - e^k_j - \xi_k\right) = 0 \label{KKT_5}.
\end{align}

From condition (\ref{KKT_1}) follows that

\begin{equation}
	d^k = t_k\left(G^k+\nu^k\right)
	\label{agg_subgr_dk}
\end{equation}

with the \emph{aggregate subgradient}

\begin{equation}
	G^k := \sum_{j \in J^k}\alpha_j g^j \quad \in \partial m_k(x^{k+1}).
\label{agg_subgr}
\end{equation}

Rewriting condition (\ref{KKT_5}) yields the \emph{aggregate error} 

\begin{equation}
	E_k := \sum_{j \in J^k} \alpha_j e^k_j = (G^k)^{\top}d^k+f(\hat{x}^k)-m_k(x^{k+1}).
\label{agg_err}
\end{equation}

Here relation (\ref{xi}) was used to replace \(\xi_k\).

The aggregate subgradient and error are used to formulate an implementable stopping condition for the bundle algorithm. The motivation behind that becomes clear with the following lemma.

\begin{lemma}\cite[Theorem 6.68]{Geiger2002}
 Let \(X = \R^n\). Let \(\varepsilon > 0\), \(\hat{x}^k \in \R^n\) and \(g^j \in \partial f(x^j)\) for \(j \in J^k\). Then the set 

\[ \mathcal{G}_{\varepsilon}^{k} := \left\{\sum_{j \in J^k}\alpha_j g^j \mid \sum_{j \in J^k} \alpha_j  e_j \leq \varepsilon, \sum_{j \in J^k}\alpha_j = 1, \alpha_j \geq 0, j \in J^k\right\} \]

is a subset of the \(\varepsilon\)-subdifferential of \(f(\hat{x}^k)\)
\[ \mathcal{G}_{\varepsilon}^k \subseteq \partial_{\varepsilon}f(\hat{x}^k). \]

\end{lemma}

This means that in the unconstrained case \(G^k \in \partial_{E_k} f(\hat{x}^k)\). So driving \(\|G^k\|\) and \(E_k\) close to zero results in some approximate \(\varepsilon\)-optimality of the objective function.
In the constrained case the stopping condition is written as

\begin{equation*}
	\delta_k = E^k+t_k\|G^k+\nu^k\|^2 \leq \mathtt{tol}.
\label{stop_cond}
\end{equation*}

\(\delta_k\) is the same measure that is also taken for the decrease test.
The relation

\begin{align*}
	\delta_k &= E^k+t_k\|G^k+\nu^k\|^2 \\
  &= E^k-\Langle G^k,d^k\Rangle - \Langle \nu^k,d^k \Rangle \\
	&= f(\hat{x}^k)-m_k(x^{k+1}) - \Langle \nu^k,d^k \Rangle \\  
\end{align*}

where (\ref{agg_subgr}) and (\ref{agg_err}) were used, shows that the new \(\delta_k\) is only a little variation of the model decrease \(\delta^M_k\).
If the iterate \(x^{k+1}\) does not lie on the boundary of the constraint set \(X\), the vector \(\nu^k = 0\) and the expression simplifies to the one stated in (\ref{mod_dec}).

For the model update the following two conditions are assumed to be fulfilled in consecutive null steps:
\begin{align}
	&m_{k+1}(\hat{x}^k+d) \geq f(\hat{x}^{k+1})-e^{k+1}_{k+1}+\left\langle g^{k+1},d \right\rangle \label{model_update_1} \\
	&m_{k+1}(\hat{x}^k+d) \geq a_k(\hat{x}^k+d) 
	\label{model_update_2}
\end{align}

The first condition means, that the newly computed information is always put into the bundle. The second one is important when updating the bundle index set \(J^k\).
It holds trivially if no or only inactive information \(j \text{ with }\alpha_j = 0\) is removed \cite{Hare2016}.
It is also always satisfied if the aggregate linearization \(a_k\) itself is added to the bundle. In this case active information can be removed without violating the condition. This is the key idea of Kiwiel's aggregation technique and ensures that the set \( \{j \in J^k| \alpha_j >0 \} \) can be bounded.

An issue of bundle methods is that in spite of the possibility to delete inactive information the bundle can still become very big. Kiwiel therefore proposed a totally different use of the aggregate objects in \cite{Kiwiel1986}. The aggregate subgradient can be used to build the \emph{aggregate linearization}

\begin{equation*}
	a_k(\hat{x}^k+d):=m_k(x^{k+1})+\langle G^k,d-d^k \rangle.
\label{agg_lin}
\end{equation*}

This function can be used to avoid memory overflow as it compresses the information of all bundle elements into one affine plane. Adding the function \(a_k\) to the cutting plane model preserves the assumptions (\ref{model_update_1}) and (\ref{model_update_2}) put on the model and can therefore be used instead of or in combination with the usual cutting planes.

This can however impair the speed of convergence if the bundle is kept too small and provides hence less information about the objective function \cite{Oliveira2014a}.

We have now all the ingredients so that the following basic bundle algorithm can be stated: 

%\textcolor{blue}{algorithm}

\begin{minipage}\linewidth

\vspace{1em}

\hrule  \vspace{0.4ex} \hrule
\vspace{1ex}
\textbf{Basic Bundle Method}
\vspace{1ex}
\hrule
\vspace{1ex}
Select a descent parameter \( m \in (0,1)\) and a stopping tolerance \( \mathtt{tol} \geq 0\). Choose a starting point \(x^1 \in \R^n\) and compute \(f(x^1)\) and \(g^1\). Set the initial index set \(J_1:=\{1\}\) and the initial stability center to \(\hat{x}^1 := x^1\), \(f(\hat{x}^1) = f(x^1)\) and select \(t_1 > 0\).
\end{minipage}

For \(k = 1,2,3  \dotsc \)   

\begin{enumerate}
	\item Calculate
	\[ d^k = \argmin_{d \in \R^n}{m_k(\hat{x}^k+d) + \mathtt{i}_X(\hat{x}^k+d) +\frac{1}{2t_k}\|d\|^2}  \]
	and the corresponding Lagrange multiplier \(\alpha_j^k,~ j \in J_k\).
	\item Set
		\begin{equation*}
			\begin{split}
				G^k = \sum_{j \in J_k}{\alpha^k_j g^k_j}, \quad E_k = \sum_{j \in J_k}{\alpha^k_j e^k_j}, \quad \text{and} \quad  \delta_k =  E_k + t_k\|G^k+\nu^k\|^2
			\end{split}
		\end{equation*}
		If \(\delta_k \leq \mathtt{tol} \rightarrow \) STOP.
	\item Set \( x^{k+1} = \hat{x}^k + d^k \).
	\item Compute \(f(x^{k+1}),~ g^{k+1}\). \\
	If 
	\[f(x^{k+1}) \leq f(\hat{x}^k) - m\delta_k \quad \rightarrow  \text{serious step}. \]
	Set \(\hat{x}^{k+1} = x^{k+1}, ~f(\hat{x}^{k+1}) = f(x^{k+1})\) and select a suitable \(t_{k+1} > 0\). \\
	Otherwise
	\[\rightarrow\text{ nullstep.} \] \\
	Set \(\hat{x}^{k+1} = \hat{x}^k, ~f(\hat{x}^{k+1})=f(x^{k+1})\) and choose \(t_{k+1} > 0\) in a suitable way.
	\item Select the new bundle index set \(J_{k+1}\), calculate \(e_j^{k+1}\) for \(j \in J_{k+1}\)	and update the model \(m_k\).
\end{enumerate}
\vspace{1ex}
\hrule

\vspace{1.5em}

In steps 4 and 5 of the algorithm it is not specified how to update the parameter \(t_k\), the index set \(J^k\) and the model \(m_k\).
For the convergence proof it is only necessary that \(\liminf_{k \to \infty} t_k > 0\) and that conditions (\ref{model_update_1}) and (\ref{model_update_2}) are fulfilled.

In practice the choice of \(t_k\) can be realized by taking

\begin{equation}
	t_{k+1} = \kappa_{+}t_k, \quad \kappa_{+} > 1
	\label{t_plus}
\end{equation}

at every serious step and 

\begin{equation}
	t_{k+1} = \max\{\kappa_{-}t_k, t_{min}\}, \quad \kappa_{-} < 1 \text{ and } t_{min} > 0
	\label{t_minus}
\end{equation}

at every null step.
The idea behind this management of \(t_k\) is taken from the trust region method: If the computed iterate was good, the model is assumed to be reliable in a bigger area around this serious iterate so bigger step sizes are allowed. If a null step was taken, the model seems to be too inaccurate far from the current serious point. Then smaller step sizes are used.
A more sophisticated version of this kind of step size management is also used by Noll et al. in \cite{Noll2012} and \cite{Noll2013}. The trust region idea was very much exploited by Schramm and Zowe in \cite{Schramm1992a}.
In the case \(X = R^n\) the sequence \(\{\hat{x}^k\}\) can be unbounded. In this case bounding \(t_k < \infty ~ \forall k\) preserves the convergence proof \cite[Theorem 3.2.2]{Hiriart-Urruty1993}.

In general it can be shown that if \(f\) posses global minima and the basic bundle algorithm generates the sequence \(\{\hat{x}^k\}\) this sequence converges to a minimizer of problem (\ref{min_prob_basic}) (c.f \cite{Hiriart-Urruty1993}).



%updates of the steplength \(t_k\), the index set \(J_k\) and the model update are are only given in a very general form. 
%The ``suitable'' choice of \(t_k\) will be discussed more closely in the convergence analysis of \textcolor{red}{decide which method; say that \(t_k > 0 \forall k...\)}. \\
%\textcolor{red}{Comment on \(J_k\) update \(\rightarrow\) depends on what is included in thesis.} \\
%For the choice of the new index set \(J_{k+1}\) different aggregation methods to keep the memory size controllable are available. The most easy and intuitive one is to just take those parts of the model function, that are actually active in the current iteration. This is done in this basic version of the method. \\

%\section{Variations of The Bundle Method}

After their discovery in 1975 bundle methods became very successful. Only a few years after that they were generalized to be used also with nonconvex objective functions. Early works, that contain fundamental ideas still used for these algorithms are \cite{Mifflin1982} and \cite{Kiwiel1985}. It then took over 25 years that bundle methods were again generalized to the use of inexact information. The first works on this subject being \cite{Hintermueller2001,Kiwiel2006} and \cite{Solodov2003}.

This section of the thesis shortly presents the key ideas of those two kinds of generalizations and different types of bundle methods that realize them.
This is first done for the case of convex objective functions with inexact function value and/or subgradient information and then for nonconvex objective functions. 


\subsection{Convex Bundle Methods with Inexact Information}

We focus here on convex bundle methods with inexact information. The reason for this is that there is a fundamental difference in treating inexactness between methods that assume convex and those that assume nonconvex objective functions.
When dealing with nonconvex objective functions inexactness is treated as some additional nonconvexity. This is not possible when assuming convex objective functions and the convexity property is to be exploited for better convergence results.
A throughout study on this subject including a synthetic convergence theory is done in \cite{Oliveira2014}. Here the most important aspects \textcolor{red}{in view of the following chapters are reviewed.}

\subsubsection{Degrees of Inexactness}

2 a) What information is inexact (only subgradients/both...) --> what do you gain from this?


2 b) What kind of assumptions are on the inexactness? (asymptotic, only over- /underestimation?)exactness only for serious steps,accuracy controllable or not

\textcolor{red}{Different ``degrees'' of inexactness: inexact subgradients; also function values (only subgradients easier???); asymptotically ecact; , not at null steps;  --> throughout study in in depth paper.}

1. How is generally dealt with inexact information in the algorithms \\

\textcolor{red}{possible simplifications of the algorithm} \\
Convergence for inexact convex functions:
\begin{itemize}
	\item states in paper \cite{Hare2016} (p. 14) that for convex functions error of \(\bar{\sigma}\) instead of \(2\bar{\sigma}\) possible (and for lower models; see depth paper?)
\end{itemize}
%\begin{itemize}
	%\item recognized: fundamentally different? approach for convex and nonconvex functions (at least in algorithm) \\
	%convex: ``deal'' with inexactness; extra steps... \\
	%nonconvex: generally no difference in algorithm (but for example line search not possible --> only no change, if algorithm was suitable before)
	%\item nonconvex algorithms: inexactness is seen as some kind of nonconvexity --> for function values clear, for subgradients???
%\end{itemize}

%\textcolor{blue}{seems to be the same:
%\begin{align}
	%& \|g_a-g\| \leq \theta \\
	%\Leftrightarrow \quad & g_a \in \partial f + B_\theta(0) \\
	%\Leftrightarrow \quad & g_a \in \partial_{\varepsilon}f, \quad \theta \leq \varepsilon^2
%\end{align}
%Last implication only for convex functions because \(\varepsilon\)-subdifferential otherwise not defined.
%See also papers from ``Chinese-search''} \\

\subsection{Nonconvex Bundle Methods with Exact Information}

There are different approaches for handling nonconvexity of the objective function in bundle methods.
As the nonnegativity property of the linearization errors \(e_j^k\) is crucial for the convergence proof of convex bundle methods an early idea was forcing the errors to be so by different downshifting strategies. A very common one is using the \emph{subgradient locality measure} \cite{Kiwiel1986, Mifflin1982}. Here the linearization error is essentially replaced by the nonnegative number

\begin{equation}
	\tilde{e}_j^k := \max_{j \in J_k} \{|e_j^k|,\gamma \|\hat{x}^k-x^j\|^2\}
\label{subgr_loc_measure}
\end{equation}

or a variation of this expression.

The expression gradient locality measure comes from the dual point of view, where the aggregate linearization error provides a measure for the 'distance' of the calculated \(\varepsilon\)-subgradient to the objective function.

Methods that use downshifting for building the model function are often endowed with a line search to provide sufficient decrease of the objective function. For the linesearch to terminate finitely, usually semismoothness of the objective function is needed.

Instead of using line search it is also possible to do \emph{proximity control}. This means that the step size parameter \(t_k\) is managed in a smart way to ensure the right amount of decrease in the objective function. This method is very helpful in the case of nonconvex objective functions with inexact information as it is predominantly considered in this thesis.

As inexactness can be seen as a kind of slight nonconvexity one could be tempted to think that nonconvex bundle methods are destined to be extended to the inexact case. Indeed, the two existing algorithms \cite{Hare2016,Noll2013} that deal with both nonconvexity and inexactness are both extensions of a nonsmooth bundle method.
This is however seldom possible for algorithms that employ a line search because for functions with inexact information convergence of this subroutine cannot be proven.

To this end proximity control seems to be a very promising strategy. It is used in many different variations in \cite{Apkarian2008, Lewis2015, Noll2005, Noll2010, Noll2012, Schramm1992}

In the beginning bundle methods were mostly explored from the  dual point of view. Newer concepts focus also on the primal version of the method. This invokes for example having different model functions for the subproblem.

In \cite{Fuduli2004, Fuduli2004a} the difference function 

\begin{equation}
	h(d):= f(x^j +d) -f(x^j) \quad j \in J_k
\label{diff_fun}
\end{equation}

is approximated to find a descent direction of \(f\).
The negative linearization errors are addressed by having two different bundles. One containing the indices with nonnegative linearization errors and one containing the other ones. From these two bundles two cutting plane approximations can be constructed which provide the bases for the calculation of new iterates.

In \cite{Noll2012} Noll et al. follow an approach of approximating a local model of the objective function. The model can be seen as a nonsmooth generalization of the Taylor expansion and looks the following:

\begin{equation}
	\Phi(y,x) = \phi(y,x)+\frac{1}{2}(y-x)^{\top}Q(x)(y-x)
\label{quad_mod}
\end{equation}

The so called \emph{first order model} \(\phi(.,x)\) is convex but possibly nonsmooth and can be approximated by cutting planes. The \emph{second order part} is a quadratic but not necessarily convex. The algorithm then proceeds a lot in the lines of a general bundle algorithm.
Instead of a line search is uses proximity control to ensure convergence.

Generally for all of this methods convergence to a stationary point is established under the assumptions of a locally Lipschitz objective function and bounded level sets \(\{x \in \R^n | f(x) \leq f(\hat{x}^1)\}\).
If the method uses a line search additionally semismoothness of the objective function is needed.

In \cite{Noll2013} the second order approach of \cite{Noll2012} is extended to functions with inexact information.
As far a we know this is the only other bundle method that can deal with nonconvexity and inexactness in both the function value and subgradient. It inspires the variable metric variation of the method used by Hare et al. in \cite{Hare2016} that is presented in section \ref{variable_metric} of this thesis.

To conclude this section we can say: At the moment there exist two fundamentally different approaches to tackle inexactness in various bundle methods depending on if the method is developed 	for convex or nonconvex objective functions.
In the nonconvex case inexactness is only considered in the paper by Hare, Sagastiz{\`{a}}bal and Sodolov \cite{Hare2016} presented above and Noll \cite{Noll2013}. In these cases the inexactness can be seen as an ``additional nonconvexity''. In practice this means that the algorithm can be taken from the nonconvex case with no or only minor changes.
In case of convex objective functions changes in the algorithm are more involved. The reason for this is that generally stronger convergence results are possible with inexactness in the convex case than in the nonconvex case. This means however, that the inexactness cannot be incorporated as easily into the algorithm.

%%\documentclass[12pt, a4paper]{scrartcl}
%\usepackage{fullpage}
%\usepackage[latin1]{inputenc} %Zeichensatzkodierung (auch Sonderzeichen -> �,�,�)
%\usepackage[english, ngerman]{babel} %Sprache -> Silbentrennung, automatisch generierte Texte auf deutsch; ngerman -> NEUE deutsche Rechtschreibung
%\usepackage[T1]{fontenc} %Sonderzeichen allg
%\usepackage{lmodern} %PDF-optimierte Schrift
%
%\usepackage{amsmath, amssymb, amsbsy, amsthm} %Mathepakete ; amsthm: Theorem-Umgebung; muss nach amsmath eingebunden werden
%\usepackage{eqnarray}
%\usepackage[x11names]{xcolor} % muss vor "`mcode"' da dieses bereits intern das color-Paket l�dt und es dann einmal mit und einmal ohne Option geladen wird, was nicht geht 
%\usepackage{graphicx} % Grafiken einf�gen
%%\usepackage{listings} % zum Einbinden von Code, ist bereits in "`mcode"' enthalten
%\usepackage{enumerate}
%\usepackage{tabularx}
%\usepackage{here}
%\usepackage{bibgerm} % F�r das Literaturverzeichnis
%% \usepackage{extarrows} % Pfeile mit Beschriftung, deren L�nge automatisch angepasst wird
%\usepackage{algorithm2e} % Pseudocode
%\usepackage{algpseudocode}
%\usepackage{bm}
%
%\usepackage{setspace}
%% um Zeilenabstand zu setzen
%
%%\usepackage{bibgerm} % F�r das Literaturverzeichnis
%
%
%
%%Definition der gew�nschten Theorem-Stile:
%\newtheorem{theorem}{Theorem}[section]
%\newtheorem{lemma}[theorem]{Lemma} %hier [section] nicht noch einmal wiederholen, da durch [satz] die Nummerierung schon genauo wie in "Satz" definiert
%\newtheorem{corollary}[theorem]{Corrolary}
%
%\theoremstyle{definition}
%\newtheorem{definition}[theorem]{Definition}
%\newtheorem{convention}[theorem]{Convention}
%\newtheorem{example}[theorem]{Example}
%
%\newtheoremstyle{my_remark}
%{}{}{}{}{\itshape}{:}{.5em}{}
%
%\theoremstyle{my_remark}
%\newtheorem*{remark}{Remark}
%
%%\newenvironment{proof}{\begin{proof}[Proof:]}{\end{proof}}
%
%
%\newcommand{\R}{\mathbb{R}}
%\newcommand{\N}{\mathbb{N}}
%\newcommand{\blangle}{\big\langle}
%\newcommand{\brangle}{\big\rangle}
%\DeclareMathOperator*{\argmin}{arg\,min}
%\newcommand{\erl}{\textcolor{Green3}{ - erledigt}}
%
%\allowdisplaybreaks[0] % Umgebungen wie align d�rfen Seitenumbr�che enthalten
%
%%\setlength{\voffset}{-28.4mm}
%%\setlength{\hoffset}{-1in}
%%\setlength{\topmargin}{20mm}
%%\setlength{\oddsidemargin}{25mm}
%%\setlength{\evensidemargin}{25mm}
%%\setlength{\textwidth}{160mm}
%%
%\setlength{\parindent}{0pt}
%\setlength{\parskip}{1ex}
%%
%%\setlength{\textheight}{235mm}
%%\setlength{\footskip}{20mm}
%%\setlength{\headsep}{50pt}
%%\setlength{\headheight}{0pt}
%
%
%
%\begin{document}
%
%
%\setcounter{tocdepth}{2} % Im Inhaltsverzeichnis werden nur section und subsection ber�cksichtigt, keine tieferen Ebenen
%\thispagestyle{empty} %keine seitenzahlen?
%\tableofcontents
%\newpage
%
%
\section{\textcolor{red}{Preliminaries}}

\textcolor{red}{Check if requirements on functions are stated and defined.}\\

Throughout this thesis I consider the optimization Problem 
\begin{equation}
\min_{x} f(x), \quad x \in X \subseteq \R^n
\label{min_prob}
\end{equation}

where \(f\) is a possibly nonsmooth function. %\textcolor{red}{For the biggest part of this thesis and if not otherwise stated \(X = \R^n\).} \\
%\textcolor{red}{Also write something about inexactness? specify \(X\) more precisely? Convex?} \\
When it comes to nonsmooth objective functions the  derivative based framework of nonlinear optimization methods does not work any more. Therefore the most important definitions and results needed when working with nonsmooth functions are stated in this section. \\
\textcolor{red}{Just definition, lemma, theorem or a bit explanation around it?} \\
\textbf{\textcolor{red}{See if requirements in definitions and theorems meet what is needed/provided later.}} \\

\begin{definition} % as in \cite{Clarke1990}
		The function \(f: \R^n \to \R\) is called \emph{Lipschitz near \(x \in \R^n\)} if there exist \(C > 0\) and \(\varepsilon > 0\) such that 
	\[ |f(y_2) - f(y_1)| \leq C\|y_2 -y_1\| \quad \forall y_1, y_2 \in \mathbf{B}_{\varepsilon}(x). \]
\end{definition}

\begin{definition} % as in \cite{Mifflin1977}
	The function \(f: \R^n \to \R\) is called \emph{locally Lipschitz} if it is Lipschitz on each bounded subset \(B \subseteq \R^n\)
	\[ |f(y) - f(x)| \leq C \|y-x\| \quad \forall x,y \in B, \quad C >0. \]
\end{definition}

\begin{definition} % as in \cite{Hiriart-Urruty1996}
	The \emph{directional derivative} of \(f\) at \(x\) in direction \(d\) is 
	\[ f'(x,d) := \lim_{\lambda \downarrow 0}\frac{f(x+\lambda d) - f(x)}{\lambda}. \]
\end{definition}

\begin{definition}
	Let \(f\) convex. The \emph{subdifferential} \(\partial f(x)\) of \(f\) at \(x\) is the nonempty compact convex set
	\[ \partial f(x) = \{g \in \R^n | f'(x,d) \geq \langle g,d\rangle \forall d \in \R^n \}. \]
\end{definition}

\begin{definition} % as in \cite{Clarke1990}
	The \emph{generalized directional derivative} of \(f\) at \(x\) in direction \(d\) is given by
	\[ f^{\circ}(x,d):= \limsup_{\substack{y \to x \\ \lambda \downarrow 0}} \frac{f(y+\lambda d)-f(y)}{\lambda}. \]
\end{definition}

\begin{definition} % as in \cite{Clarke1990}
	The \emph{generalized gradient} of \(f\) at \(x\) is a nonempty convex compact set \(\partial f(x)\) given by
	\[ \partial f(x) := \{g \in \R^n | f^{\circ}(x,d) \geq \langle g,d\rangle \forall d \in \R^n \}. \]
\end{definition}

If \(f\) is a convex function the generalized gradient coincides with the subdifferential \(\partial f\) of \(f\) \cite{Clarke1990}.

\section{Bundle Methods}

\textcolor{Green3}{introduction}

%\textcolor{red}{\(\rightarrow\) introduction sentences: main focus on algorithm  in paper; therefore at first introduce simplest form of bundle methods.}
When bundle methods were first introduced in 1975 by Claude Lemar�chal and Philip Wolfe they were developed to minimize a convex (possibly nonsmooth) function \(f\) for which at least one subgradient at any point \(x\) can be computed \cite{Mifflin2012}. \\
To provide an easier understanding of the proximal bundle method in \cite{Hare2016} and stress the most important ideas of how to deal with nonconvexity and inexactness first a basic bundle method is shown here. \\
Bundle methods can be interpreted in two different ways: From the dual point of view one tries to approximate the \(\varepsilon\)-subdifferential to finally ensure first order optimality conditions. The primal point of view interprets the bundle method as a stabilized form of the cutting plane method where the objective function is modeled by tangent hyperplanes \cite{Hare2010}. I focus here on the primal approach. \\
%\textcolor{red}{In the next two sections the function \(f\) is assumed to be convex.} \\

%\textcolor{blue}{notation, definitions}\\
%\textcolor{red}{already done in previous preliminaries chapter?}

\subsection{A basic bundle method}

\textcolor{red}{This section gives a short summery of the derivations and results of chapter XV in \cite{Hiriart-Urruty1993} where a primal bundle method is derived as a stabilized version of the cutting plane method. If not otherwise indicated the results in this section are therefore taken from \cite{Hiriart-Urruty1993}.} 

\textcolor{red}{The optimization problem considered in this section is 
\begin{equation}
  \min_{x} f(x) \quad \text{s.t.} \quad x \in X
\label{min_prob_basic}
\end{equation}
with the convex function \(f\) and the closed and convex set \(X \subseteq \R^n\).}

%\textcolor{red}{Define Problem again?? Incorporate ``set-constraint'' by writing \(h(x):= f(x)+\mathbb{I}_X\). \(\rightarrow\) later???}

\textcolor{blue}{explanation}

\subsubsection{Derivation of the bundle method}

The geometric idea of the cutting plane method is to build a piecewise linear model of the objective function \(f\) that can be minimized more easily than the original objective function. \\
This model is built from a \emph{bundle} of information that is gathered in the previous iterations. \\
In the \(k\)'th iteration, the bundle consists of the previous iterates \(x^j\), the respective function values \(f(x^j)\) and a subgradient at each point \(g^j \in \partial f(x^j)\) for all indices \(j\) in the index set \(J_k\). From each of these triples, one can construct a linear function 

\begin{equation}
	l_j(x) = f(x^j) + (g^j)^{\top}(x-x^j)
\label{lin_fun}
\end{equation}  

with \(f(x^j) = l_j(x^j)\) and due to convexity \(f(x) \geq l_j(x), ~ x \in X\).\\ % x \in \R^n??
One can now model the objective function \(f\) by the piecewise linear function

\begin{equation}
	m_k(x) = \max_{j \in J_k} l_j(x)
\label{cp_model}
\end{equation}

and find a new iterate \(x^{k+1}\) by solving the subproblem

\begin{equation}
	\min_{x} m_k(x) \quad \text{s.t.} \quad x \in X.
\label{cp_model_fun}
\end{equation}

This subproblem should of course be easier to solve than the original task. A question that depends a lot on the structure of \(X\). If \(X = \R^n\) or a polyhedron, the problem can be solved easily. Still there are some major drawbacks to the idea. For example if \(X = \R^n\) the solution of the subproblem in the first iteration is always \(-\infty\). \\
In general one can say that the subproblem does not necessarily have to have a solution. \\
To tackle this problem a penalty term is introduced to the subproblem:

\begin{equation}
	\tilde{m}_k(x) =  m_k(x) + \frac{1}{2 t}\|x-x^k\|^2 \quad \text{s.t.} \quad x \in X
\label{model_fun}
\end{equation}

This new subproblem is strongly convex and has therefore always a unique solution.\\
\textcolor{red}{how much explanation here?}
% \max_{j \in J_k} l_j(\hat{x}^k + d)

%\textcolor{red}{Some nice sentences to explain the term a little bit more and to lead over to the next paragraph.\\ 
%To understand the deeper motivation of this term see \cite{Hiriart-Urruty1993}. For this introduction it suffices to see that due to the regularization term the subproblem is now strongly convex and therefore always uniquely solvable.} \\


The second major step towards the bundle algorithm is the introduction of a so called \emph{stability center} or  \emph{serious point} \(\hat{x}^k\). It is the iterate that yields the ``best'' approximation of the optimal point up to the \(k\)'th iteration (\textcolor{red}{not necessarily the best function value though}).\\
The updating technique for \(\hat{x}^k\) is crucial for the convergence of the method: If the next iterate yields a decrease of \(f\) that is ``big enough'', namely bigger than a fraction of the decrease suggested by the model function for this iterate, the stability center is moved to that iterate. If this is not the case, the stability center remains unchanged. \\
In practice this looks the following: \\
Define first the \emph{nominal decrease} \(\delta_k\) which is the decrease of the model for the new iterate \(x^{k+1}\) compared to the function value at the current stability center \(\hat{x}^k\).

\begin{equation}
	\delta^{k} = f(\hat{x}^k) - \tilde{m}_k(x^{k+1}) +a_k \geq 0
\label{nom_dec}
\end{equation}

\textcolor{red}{The nominal decrease is in fact stated a little differently for different versions of the bundle algorithm, this is why I added the constant \(a_k \in \R\) here for generalization. In practice the difference between the decreases is not influencing the algorithm as \(\delta_k\) is weighted by the constant \(m \in (0,1)\) for the descent test which compensates \(a_k\).} 

If the actual decrease of the objective function is bigger than a fraction of the nominal decrease 
\[ f(\hat{x}^k) - f(x^{k+1}) \geq m \delta_k, \quad m \in (0,1) \]
set the stability center to \( \hat{x}^{k+1} = x^{k+1}\). This is called a \emph{serious} or \emph{descent step}. \\
If this is not the case a \emph{null step} is executed and the serous iterate remains the same \(\hat{x}^{k+1} = \hat{x}^k\).\\

The subproblem can be rewritten as a smooth optimization problem. For convenience rewrite the affine functions \(l_j\) with respect to the stability center \(\hat{x}^k\).

%\textcolor{red}{citation for this???!!!}

\begin{align}
	l_j(x) &= f(x^j) + {g^j}^{\top} (x - x^j) \\
	&= f(\hat{x}^k)+{g^j}^{\top}(x-\hat{x}^k)-(f(\hat{x}^k)-f(x^j) + {g^j}^{\top}(x^j-\hat{x}^k)) \\
	&= f(\hat{x}^k) + {g^j}^{\top}(x-\hat{x}^k) - e^k_j
\end{align}	

where
\begin{equation}
	e^k_j = f(\hat{x}^k)-f(x^j) + {g^j}^{\top}(x^j-\hat{x}^k) \geq 0 \quad \forall j \in J_k
\label{lin_err}
\end{equation}

is the \emph{linearization error}. The nonnegativity property is essential for the convergence theory and will also be of interest when moving on to the case of nonconvex and inexact objective functions. \\


The subproblem can now be written as

\begin{align}
	& \min_{\hat{x}^k+d \in X} \tilde{m}_k(d) = f(\hat{x}^k) + \max_{j \in J_k} \{{g^j}^{\top}d - e^k_j\} + \frac{1}{2t_k}\|d\|^2
	\label{sub_prob_long}\\
	\Leftrightarrow \quad &\min_{\hat{x}^k+d \in X} \xi + \frac{1}{2t_k}\|d\|^2 \quad \text{s.t.} \quad f(\hat{x}^k)+{g^j}^{\top}d - e^k_j - \xi \leq 0, \quad j \in J_k
\label{sub_prob_short}
\end{align}


%\textcolor{red}{where the constant term \(f(\hat{x}^k)\) was discarded for the sake of simplicity.} \\
If \(X\) is a polyhedron this is a quadratic optimization problem that can be solved using standard methods of nonlinear optimization. The pair \((\xi_k, d^k)\) solves (\ref{sub_prob_short}) if and only if \(d^k\) solves the original subproblem (\ref{sub_prob_long}) and \(\xi_k=f(\hat{x}^k)+\max_{j \in J_k}{g^j}^{\top}d^k - e_j^k\). The new iterate is then given by \(x^{k+1} = \hat{x}^k + d^k\). \\

%\begin{remark}
	%\textcolor{red}{Setting \(\check{f}(x) = f(x)+ \mathbb{I}_{X}(x)\) the above optimization problem is ...\\
	%The \emph{proximal point mapping} or \emph{prox-operator}
	%\begin{equation}
		%prox_{t,f}(x) = \argmin_{y}\left\{\check{f}(y) + \frac{1}{2t}\|x-y\|^2\right\}, \quad t > 0
	%\label{prox_op}
	%\end{equation}}
	%\textcolor{red}{source??? This special form of the subproblems gives the proximal bundle method its name and will occur again later??? }
%\end{remark}
%


\subsubsection{Aggregate objects}

The constraint \(\hat{x}^k+d \in X\) can also be incorporated directly in the objective function by using the indicator function

\[ \mathbb{I}_X(x) = \left\{ \begin{array}{cl} 0, & \text{if } x \in X \\ +\infty, & \text{if } x \notin X \end{array} \right. .\]

The subproblem then writes as
\begin{equation}
	\min_{\hat{x}^k+d \in R^n} \xi + \mathbb{I}_X + \frac{1}{2t_k}\|d\|^2 \quad \text{s.t.} \quad {g^j}^{\top}d - e^k_j - \xi \leq 0, \quad j \in J_k
\label{sub_prob_fin}
\end{equation}

%\textcolor{red}{Some introduction how this and the aggregate error expression relate to each other. Why it is in this case easier to write the model in the nonsmooth form...}

%\textcolor{red}{Lemma XI 3.1.1 \(\partial g = \partial f + \partial \mathbb{I}_X\) for \(g = f + \mathbb{I}_X\).}

One gets the following results about the step \(d^k\) of the subproblem:

\begin{lemma} %Lemma XV 3.1.1 \cite{Hiriart-Urruty1993}
	The optimization problem (\ref{sub_prob_fin}) has for \(t_k > 0\) a unique solution given by
	\begin{equation}
		d^k = -t_k(G^k + \nu^k), \quad G^k \in \partial m_k(d^k), \quad \nu^k \in \partial  \mathbb{I}_X.
	\label{sub_prob_sol}
	\end{equation}
\end{lemma}

Furthermore
\begin{equation}
	m_k(\hat{x}^k+d) \geq f(\hat{x}^k) + {G^k}^{\top}d - E_k \quad \forall d \in \R^n
\label{m_inequ}
\end{equation}

%\textcolor{red}{inequality because of aggregation technique. Is sharp when cutting plane model is used? source?}

where 
\begin{equation}
	E_k := f(\hat{x}^k) - m_k(x^{k+1}) + {G^k}^{\top}d^k.
\label{agg_err_1}
\end{equation}

%\textcolor{red}{Comment on the inequality missing}

The quantities \(G^k\) and \(E^k\) are the \emph{aggregate subgradient} and the \emph{aggregate error}.

\textcolor{red}{Explain aggregation process in more detail}

From the Karush-Kuhn-Tucker conditions (KKT-conditions) one can see that in the optimum there exist Lagrange or \emph{simplicial multiplier} \(\alpha_j^k, ~j \in J_k\) such that

\begin{equation}
	\alpha_j^k \geq 0, \quad \sum_{j \in J_k}{\alpha_j^k} = 1
\label{alphas}
\end{equation}

 %\textcolor{red}{by rewriting and so on... one can see that the above expressions are in fact }

From the dual problem one obtains that the aggregate subgradient and error can also be expressed as

\begin{equation}
	E_k = \sum_{j \in J_k} \alpha_j^k e_j^k \qquad \text{and}\qquad G^k = \sum_{j \in J_k} \alpha_j^k g^j.
\label{agg_obj}
\end{equation}

%\textcolor{red}{Finally use Lemma ??? in \cite{Hiriart-Urruty1993}} 
%\[ m_k(x^{k+1}) = f(\hat{x}^k) - E_k - t_k\|G^k\|^2 \]
%to reformulate the nominal decrease \(\delta_k\):
%\[ \delta_k =  f(\hat{x}^k) - m_k(x^{k+1}) - \frac{1}{2}t_k\|G^k\|^2 = E_k + \frac{1}{2}t_k\|G^k\|^2\]
The nominal decrease in this case is defined as:\\
%\textcolor{red}{noch mal anschauen}
\begin{equation}
	\delta_k := E_k + t_k\|G^k+\nu^k\|^2 = f(\hat{x}^k) - m_k(x^{k+1}) - {\nu^k}^{\top}d^k
\label{nom_dec_1}
\end{equation}

%\textcolor{red}{In practice the different definition of the decreases makes no difference because of the weighting with the descent parameter \(m\).}

The following basic bundle algorithm can now be stated: 

%\textcolor{red}{Reformluate equations, model function \\
%introduce aggregate expressions \\
%say something to \(J\)-update, say something to \(t\)-update \\
%see if all abbreviations (\(f_j, g^j, ...\)) are introduced \\
%introduce prox-operator and proximal points}

\textcolor{blue}{algorithm}

\vspace{1em}

\hrule  \vspace{0.4ex} \hrule
\vspace{1ex}
\textbf{Basic bundle method}
\vspace{1ex}
\hrule
\vspace{1ex}
Select descent parameter \( m \in (0,1)\) and a stopping tolerance \( \mathtt{tol} \geq 0\). Choose a starting point \(x^1 \in \R^n\) and compute \(f(x^1)\) and \(g^1\). Set the initial index set \(J_1:=\{1\}\) and the initial stability center to \(\hat{x}^1 := x^1\), \(f(\hat{x}^1) = f(x^1)\) and select \(t_1 > 0\).

For \(k = 1,2,3  \dotsc \)   

\begin{enumerate}
	\item Calculate
	\[ d^k = \argmin_{d \in \R^n}{m_k(\hat{x}^k+d) + \mathbb{I}_X +\frac{1}{2t_k}\|d\|^2}  \]
	and the corresponding Lagrange multiplier \(\alpha_j^k,~ j \in J_k\).
	%\textcolor{red}{say how model \(m_k\) looks here. include \(\ \mathbb{I}_X\)}
	\item Set
		\begin{equation*}
			\begin{split}
				G^k = \sum_{j \in J_k}{\alpha^k_j g^k_j}, \quad E_k = \sum_{j \in J_k}{\alpha^k_j e^k_j}, \quad \text{and} \quad  \delta_k =  E_k + t_k\|G^k+\nu^k\|^2
			\end{split}
		\end{equation*}
		If \(\delta_k \leq \mathtt{tol} \rightarrow \) STOP.
	\item Set \( x^{k+1} = \hat{x}^k + d^k \).
	\item Compute \(f(x^{k+1}),~ g^{k+1}\). \\
	If 
	\[f^{k+1} \leq \hat{f}^k - m\delta_k \quad \rightarrow  \text{serious step}. \]
	Set \(\hat{x}^{k+1} = x^{k+1}, ~f(\hat{x}^{k+1}) = f(x^{k+1})\) and select suitable \(t_{k+1} > 0\). \\
	Otherwise \(\rightarrow\) nullstep. \\
	Set \(\hat{x}^{k+1} = \hat{x}^k, ~f(\hat{x}^{k+1})=f(x^{k+1})\) and choose \(t_{k+1}\) in a suitable way.
	\item Select new bundle index set \(J_{k+1} = \{j \in J_k| \alpha_j^{k+1} \neq 0\} \cap {k+1}\), calculate \(e_j\) for \(j \in J_{k+1}\)	and update the model \(m_k\).
\end{enumerate}
\vspace{1ex}
\hrule

\vspace{1.5em}
In steps 4 and 5 of the algorithm the updates of the steplength \(t_k\) and the index set \(J_k\) are are only given in a very general form.\\
The ``suitable'' choice of \(t_k\) will be discussed more closely in the convergence analysis of \textcolor{red}{decide which method; say that \(t_k > 0 \forall k...\)}. \\
%\textcolor{red}{Comment on \(J_k\) update \(\rightarrow\) depends on what is included in thesis.} \\
For the choice of the new index set \(J_{k+1}\) different aggregation methods to keep the memory size controllable are available. The most easy and intuitive one is to just take those parts of the model function, that are actually active in the current iteration. This is done in this basic version of the method. \\
\textcolor{red}{Refer to low memory bundling if later in thesis.}
%Instead of keeping every index in the set \(J_k\) different compression ideas exist.  \textcolor{red}{For now I therefor stick to this update. \\
%refer to later ``low memory'' thing??}

%\textcolor{red}{explanation to \(t_k\) update. \(\rightarrow\) include at which point???} 
%This simple idea has however some major drawbacks \cite{Hiriart-Urruty1996}:
%\begin{itemize}
	%\item Minimization of the cutting plane model of the objective function is not trivial. Indeed unconstrained minimization of the model is never possible in the first step, where it is just a line, unless the starting point is already a minimum.
	%\extcolor{red}{\item The convergence speed is very slow.}
%\end{itemize}
%
%\textcolor{red}{If convergence speed named here, does it have to be shown (rates)? For all algorithms??? Leave out? Argue about instability?}
%
%To address those issues a regularization is added to the cutting plane model. This ensures unique solvability of the minimization of the subproblem. \textcolor{red}{By introducing a stability center and }
\subsection{Proximal bundle method for nonconvex functions with inexact information}

\textcolor{blue}{introduction}

This section focuses on the proximal bundle method presented in \cite{Hare2016}.\\
The idea is to extend the basic bundle algorithm for nonconvex functions with both inexact function and subgradient information. \\
The key idea of the algorithm is the one already developed for \cite{Hare2010}: When dealing with nonconvex functions a very critical difference to the convex case is that the linearization errors are not necessarily nonnegative any more. To tackle this problem the errors are manipulated to enforce nonnegativity. In this case this is done my modeling not the objective function directly but a convexified version of it. \\
 % citation for that??? sources of nonconv-exact?
%\textcolor{red}{something about primal-dual view???}
\subsubsection{New subsubsection?}
\textcolor{blue}{``assumptions and notations''}

\textcolor{red}{introduce exact optimization problem that is used in this section and its properties if not already introduce in ``Preliminaries''.}

Throughout this section the optimization problem

\begin{equation}
	\min_{x} f(x) \quad \text{s.t.} \quad x \in X
\label{opt_prob_nonconv_inex}
\end{equation}

where \(f\) is locally Lipschitz is considered. \(X \subseteq \R^n\) is assumed to be a convex compact set. Both the function value as well as the subgradient can be provided in an inexact form. \\

For the function value inexactness is defined straight forwardly: If 
\begin{equation}
	\|\tilde{f} - f(x)\| \leq \sigma
\label{fun_inex}
\end{equation}

then \(\tilde{f}\) approximates the value \(f(x)\) within \(\sigma\).\\
For the subgradients inexactness is interpreted in the following way: \(\tilde{g} \in \R^n\) approximates a subgradient \(g \in \partial f(x)\) within \(\theta \geq 0\) if
\begin{equation}
	\tilde{g} \in \partial f(x) + B_{\theta}(0).
\label{subgr_inex}
\end{equation}

In the paper it is assumed that the errors are bounded although the bound does not have to be known.

\begin{equation}
	|\sigma_j| \leq \bar{\sigma} \quad \text{and} \quad 0 \leq \theta_j \leq \bar{\theta} \quad \forall j \in J_k.
\label{err_bound}
\end{equation}

In the context of inexact information it is important to make a distinction between the (unknown) exact function value and its approximation. Throughout this chapter I therefore write \(f(x)\) for the exact function value whereas the approximation will be written as \(f_j\) or \(\hat{f}_k\) for the approximation at the current stability center.

The objective function \(f:\R^n \to \R\) is assumed to be proper, (subdifferentially) regular and locally Lipschitz continuous with full domain.

\begin{definition} \cite{Rockafellar2009}
	A function \(f:\R^n \to \bar{\R} = [-\infty, + \infty]\) is called \emph{proper} if \(f(x) < \infty\) for at least one \(x\in \R^n\) and \(f(x) > \infty ~ \forall x \in \R^n\).
\end{definition}

\begin{definition} \cite{Rockafellar2009}
	\(f: \R^n \to \bar{\R}\) is called \emph{subdifferentially regular} at \(\bar{x}\) if \(f(\bar{x})\) is finite and the epigraph
	\[epi(f) := \{(x, \alpha) \in \R^n\times \R | \alpha \geq f(x)\}\]
	is Clarke regular at \(\bar{x},f(\bar{x})\).
\end{definition}

\textcolor{red}{Closed convex sets are Clarke regular, so in particular the epigraph of \textcolor{red}{lower \(\mathcal{C}^2\)-functions?}.}

\textcolor{red}{Definition semismooth for later:} \\
\begin{definition} % as in \cite{Mifflin1977}
	A function \(f: \R^n \to \R\) is called \emph{semismooth} ar \(x\in \R^n\) if \(f\) is Lipschitz near \(x\) and for each \(d \in \R^n\) and for any sequences \(\{t_k\} \subseteq \R_{+}, \{\theta^k\} \subseteq \R^n\) and \(\{g^k\} \subseteq \R^n\) such that 
	\[ \{t_k\} \downarrow 0, \quad \{\theta^k/t_k\} \to 0 \in \R^n \quad \text{and} \quad g^k \in \partial f (x+t_kd+\theta^k), \]
	the sequence \(\{\langle g^k, d \rangle\}\) has exactly one accumulation point.
\end{definition}

\begin{definition} % as in \cite{Haarala2004a}
A point \(x \in \R^n\) that satisfies \(0 \in \partial f(x)\) is called a \emph{stationary point} of \(f\).
\end{definition}

\textcolor{blue}{explanation}

A main issue both nonconvexity and inexactness entail is that the linearization errors \(e_j^k\) are not necessarily nonnegative any more.\\
So based on the results in \cite{Hare2009} not the objective function but a convexified version of it is modeled as the objective function of the subproblem.\\
%\textcolor{red}{explain locally convexified more precisely? Is it because no global convexification? different than in [18]??}\\

When looking at the subproblem formulated as in (\ref{sub_prob_long}) one can see that the new iterate \(x^{k+1}\) is in fact a \emph{proximal point} of the subproblem.\\
The \emph{proximal point mapping} or \emph{prox-operator} is defined as
\begin{equation}
	prox_{t,f}(x) = \argmin_{y}\left\{\check{f}(y) + \frac{1}{2t}\|x-y\|^2\right\}, \quad t > 0
\label{prox_op}
\end{equation}

For \(\check{f}(x) := m_(x)+\mathbb{I}_X(x)\) and \(\mu:=\frac{1}{t_k}\) this is just subproblem (\ref{sub_prob_long}) with the constraint \(x\in X\) incorporated in the objective function. Because of this special form of the subproblems primal bundle methods are also called proximal bundle methods.\\


\textcolor{red}{explain in much more detail when read about calculation of proximal points for nonconvex functions. At the moment just main ideas.}

The key idea is now to use the relation
\begin{equation}
	prox_{R = \mu+\eta, f}(x) = prox_{\mu,f+\eta/2 \cdot |\cdot - x|^2}(x).
\label{prox_relation}
\end{equation}

This means, that the proximal point of the function \(f\) for parameter \(R=\eta+\mu\) is the same as calculating the proximal point of the regularized function 

\begin{equation}
	\tilde{f}(y) = f(y) + \frac{\eta}{2}|y-x|^2
\label{conv_obj}
\end{equation}

with respect to the parameter \(\mu\). \(\eta\) is therefore called the \emph{convexification parameter} and \(\mu\) is the \emph{prox-parameter}.\\

So the function that will be modeled by the cutting plane approximation is no longer the original objective function \(f\) but the convexified version \(\tilde{f}\). 

%\textcolor{red}{say why/how... this is related to current stability center}
%\textcolor{red}{Because new function to be approximated, subgradients ``new'':}
The linear functions forming the model have therefore a tilted slope 

\begin{align}
	s^k_j &= g^j + \eta_k \left(x^j-\hat{x}^k\right).
	\label{aug_subgr}
\end{align}

%\textcolor{red}{linearization error defined just as to be nonnegative. \(\rightarrow\) any further motivation???\\
%be careful now slightly different definition of linearization error because of inexact information}

\textcolor{red}{\(\eta\) is defined to be such that the augmented linearization error is nonnegative:
\begin{equation}
	\eta_k \geq \max\left\{\max_{j \in J_k, x^j \neq \hat{x}^k}{ \frac{-2e_j^k}{\|x^j-\hat{x}^k\|^2}}, 0 \right\} + \gamma 
\label{eta}
\end{equation}}
\textcolor{red}{With the ``saveguarding parameter'' \(\gamma \geq 0\)}

\textcolor{red}{explain, why linearization errors are defined like that}

\begin{equation}
	0 \leq c^k_j := e^k_j + b_j^k, \quad \text{with} \quad \left\{ \begin{array}{l} e_j^k := \hat{f}_k - f_j - \langle g^j, \hat{x}^k-x^j\rangle	 \vspace{1ex} \\
	b_j^k := \frac{\eta_k}{2}\|x^j-\hat{x}^k\|^2 \end{array}\right.
\label{aug_lin_err}
\end{equation}

\textcolor{red}{The new model function can therefore be written as
\begin{equation}
  M_k(d) := \hat{f}_k + \max_{j \in J_k} \left\{{s^k_j}^{\top}d-c^k_j \right\}
\label{aug_model}
\end{equation}}

%\textcolor{red}{check that \(M_k\) and \(m_k\) from basic algorithm above have same form.}\\

%\textcolor{red}{Explain how the \(D\) comes into the whole thing \\
%here (also) because already said above (but why?) \\
%already introduce in basic algorithm!\\
%maybe later say something to how it is done in ``depth''-Paper???}
%
%\textcolor{red}{Some Lemma why this is so}
%
%\textcolor{red}{ \begin{equation}
  %d^k = -t_k(G^k + \nu^k), \quad \text{where} \quad \nu^k \in \partial \mathbb{I}_D(x^{k+1})
%\label{d^k}
%\end{equation}}

\textcolor{red}{The definition of the aggregate objects follows straightforward:}
\begin{align}
	S^k := \sum_{j \in J_k} \alpha_j^k s_j^k \label{aug_agg_subgr} \\
	C_k := \sum_{j \in J_k}{\alpha_j^k c_j^k} \label{aug_agg_err}
\end{align}

\textcolor{red}{explain how \(\delta_k\) is derived.}
\textcolor{red}{\begin{equation}
	\delta^k := C_k + t_k \|S^k + \nu^k\|^2
\label{nom_dec2}
\end{equation}}

\textcolor{blue}{algorithm}

\vspace{1em}

\hrule  \vspace{0.4ex} \hrule
\vspace{1ex}
\textbf{Nonconvex proximal bundle method with inexact information}
\vspace{1ex}
\hrule
\vspace{1ex}
Select parameters \( m \in (0,1), \gamma > 0 \) and a stopping tolerance \( \mathtt{tol} \geq 0\). \\
Choose a starting point \(x^1 \in \R^n\) and compute \(f_1\) and \(g^1\). Set the initial index set \(J_1:=\{1\}\) and the initial prox-center to \(\hat{x}^1 := x^1\), \(\hat{f}_1 = f_1\) and select \(t_1 > 0\).

For \(k = 1,2,3,  \dotsc \)   

\begin{enumerate}
	\item Calculate \[d^k = \arg \min_{d \in \R^n} \left\{ M_k(\hat{x}^k+d)+\mathbb{I}_X(\hat{x}^k+d)+\frac{1}{2t_k}\|d\|^2\right\}.\]
	\item Set
		\begin{align*} 
		  G^k &= \sum_{j \in J_k}{\alpha_j^k s_j^k}, \quad	\nu^k = -\frac{1}{t_k}d^k-G^k\\
			C_k &= \sum_{j \in J_k}{\alpha_j^k c_j^k} \\
	    \delta_k &=  C_k + t_k\|G^k + \nu^k\|^2
		\end{align*}
		If \(\delta_k \leq \mathtt{tol} \rightarrow \) STOP.
	\item Set \( x^{k+1} = \hat{x}^k + d^k \).
	\item Compute \(f^{k+1}, g^{k+1}\) \\
	If 
	\[f^{k+1} \leq \hat{f}^k - m\delta_k \quad \rightarrow \text{ serious step} \]
	Set \(\hat{x}^{k+1} = x^{k+1}, \hat{f}^{k+1} = f^{k+1}\) and select \(t_{k+1} > 0\). \\
	Otherwise \(\rightarrow\) nullstep \\
	Set \(\hat{x}^{k+1} = \hat{x}^k, \hat{f}^{k+1}=f^{k+1}\) and choose \(0 < t_{k+1} \leq t_k\). 	
	\item Select new bundle index set \(J_{k+1}\), keeping all active elements. Calculate 
	\[ \eta_k \geq \max{\left\{\max_{j \in J_{k+1}, x^j \neq \hat{x}^{k+1}}{\frac{-2e_j^k}{|x^j - \hat{x}^{k+1}|^2}, 0}\right\}}+\gamma  \]
	and update the model \(M^k\)
\end{enumerate}
\vspace{1ex}
\hrule

\vspace{1.5em}

\subsection{Convergence analysis}

\subsubsection{Results for objectives with exact information}

The main ideas of the algorithm are basicly the ones developed in \cite{Hare2010} for the redistributed proximal bundle method for exact nonconvex problems. \\
Setting the error bounds \(\bar{\sigma}\) and \(\bar{\theta}\) to zero results therefore in the following convergence theorem. 

\begin{theorem}
	Let the sequence \(\{\eta_k\}\) be bounded, \(\liminf_{k \to \infty }\) and the cardinality of the set \(\{j \in J_k | \alpha_j^k > 0\}\) be uniformly bounded in \(k\). \\
	Then every accumulation point of sequence of serious iterates \(\{\hat{x}^k\}\) is a stationary point of the problem.
\end{theorem}

\textcolor{red}{think last condition only interesting in inexact case.}

In the exact case boundedness of the sequence \(\{\eta_k\}\) is proven for lower-\(\mathcal{C}^2\) functions in \cite{Hare2010}. This is not possible in the inexact case, even if the objective function \(f\) is convex.

A further simplification of the method for exact information is not necessary as the method is already almost as simple as the basic bundle method for nonconvex exact functions. Additionally no new concepts needed to be introduced when doing the step from nonconvex exact problems, for which the algorithm was originally designed, to problems with inexact information.

\textcolor{red}{\begin{remark}
	I want to add here, that the simplicity of the algorithm is rather special for methods suitable for nonconvex problems. Often a linesearch algorithm has to be inserted in the nonconvex case, which is not needed here. 
\end{remark}}

\subsubsection{Nonconvex bundle methods}
There are different approaches for handling nonconvexity of the objective function in bundle methods.
As the nonnegativity property of the linearization errors \(e_j^k\) is crucial for the convergence proof of convex bundle methods an early idea was forcing the errors to be so by different downshifting strategies. A very common one is using the \emph{subgradient locality measure} \cite{Kiwiel1986, Mifflin1982}. Here the linearization error is essentially replaced by the nonnegative number

\begin{equation}
	\tilde{e}_j^k := \max_{j \in J_k} \{|e_j^k|,\gamma \|\hat{x}^k-x^j\|^2\}.
\label{subgr_loc_measure}
\end{equation}

\textcolor{red}{Remark on dual view? How subgradient locality measure measures how close subgradient is to subdifferential of \(f\)???} \\
Methods using this kind of manipulation of the model function are endowed with a line search to provide sufficient decrease of the objective function. For the linesearch to terminate finitely, semismoothness of the objective function is usually needed. \\
It can be proven that every accumulation point of the sequence of serious points \(\{\hat{x}^k\}\) is a stationary point of the objective function \(f\) under the additional assumptions that \(f\) is locally Lipschitz and the level set \(\{x \in \R^n | f(x) \leq f(\hat{x}^1)\}\) is bounded \cite{Haarala2007}.

A drawback to the method described a bove is that it is primarily supported from the dual point of view of the bundle algorithm. Newer concepts focus also on the primal point of view. This invokes for example having different model functions for the subproblem.

In \cite{Fuduli2004, Fuduli2004a} the difference function 

\begin{equation}
	h(d):= f(x^j +d) -f(x^j) \quad j \in J_k
\label{diff_fun}
\end{equation}

is approximated to find descent direction of \(f\). \\
The negative linearization errors are addressed by having two different bundles. One containing the indices with nonnegative linearization errors and one containing the other ones. From these two bundles two cutting plane approximations can be constructed which provide the bases for the calculation of the new iterate.  \\
Convergence of the method to a stationary point is proven under the assumption of \(f\) being locally Lipschitz and semismooth. 

In \cite{Noll2012} Noll et al. follow an approach of approximating a local model of the objective function. The model can be seen as a nonsmooth generalization of the Taylor expansion and looks the following:

\begin{equation}
	\Phi(y,x) = \phi(y,x)+\frac{1}{2}(y-x)^{\top}Q(x)(y-x)
\label{quad_mod}
\end{equation}

The so called \emph{first order model} \(\phi(.,x)\) is convex but possibly nonsmooth and can be approximated by cutting planes. The \emph{second order part} is a quadratic but not necessarily convex. The algorithm then proceeds a lot in the lines of a general bundle algorithm. \\
For a locally Lipschitz objective function with a bounded levelset \(\{x \in \R^n | f(x) \leq f(\hat{x}^1)\}\) convergence to a stationary point is established.

%Other papers reach the same results only needing \textcolor{red}{L-continuity and the boundedness of the level sets (in \cite{Hare2016} established by introducing the set \(D\)). But they use a different concept.} \\
%\textcolor{red}{Explain ``most common?'' concepts for dealing with nonconvexity??? see \cite{Mifflin1982}.}

%\begin{itemize}
	%\item in paper other method for calculating \(\mu = \frac{1}{t_k}\) \\
	%\textcolor{red}{how is made sure that \(\eta\) big enough???}
	%\item convergence results for nonconvex functions:?
	%\begin{itemize}
		%\item conditions on functions: \\
		%\emph{exact}: \(f\) lower-\(\mathcal{C}^2\) near the minimizers of the problem \\
		%\emph{inexact}: proper, regular, locally Lipschitz with full domain; even better: lower-\(\mathcal{C}^1\) (contains lower-\(\mathcal{C}^2\)) \(\rightarrow\) conditions more general than in exact case
		%\item convergence results: \\
		%\emph{exact}: the limit of the sequence \(\{x^{k}\}\) (which exists) or every accumulation point of the sequence \(\{\hat{x}^k\}\) is a stationary point of \(f\)\\
		%\textcolor{red}{Does this mean: \(0 \in \partial f(\bar{x}), ~\bar{x}\) being the respective limit??? incorporate set \(D\)?} \\
		%\emph{inexact}: \(0 \in \left(\partial f(\bar{x}) + \partial I_D(\bar{x}) \right) + B_{\bar{\theta}}(0)\) \\
		%if \(f\) lower-\(\mathcal{C}^1\): 
		%\[ \forall \varepsilon >0 ~\exists \rho >0: \quad f(y) \geq f(\bar{x})-(\bar{\theta}+\varepsilon)\|y-\bar{x}\|-2\bar{\sigma} \quad \forall y \in D \cap B_{\rho}(\bar{x}) \]
		%\item obvious difference: no error terms in exact case
	%\end{itemize}
%\end{itemize}

%\textbf{Results:}
%\begin{itemize}
	%%\item first result seen: leave uncertainty away
	%\item better in comparison to other paper, because wider range of functions (lower-\(\mathcal{C}^1\)).
%\end{itemize}

\subsubsection{if convex function}

Convergence for inexact convex functions:
\begin{itemize}
	\item states in paper \cite{Hare2016} (p. 14) that for convex functions error of \(\bar{\sigma}\) instead of \(2\bar{\sigma}\) possible (and for lower models; see depth paper?)
\end{itemize}


\textbf{To Do:}
\begin{itemize}
	\item \textcolor{Green4}{proof serious steps}
	\item \textcolor{Green4}{proof null steps}
	\item \textcolor{Green4}{limit of \(G^k\)}
	\item proof in book; see if possible to leave out \(D\); compare \\
	should be possible if bounded level sets assumed; check this! \\
	compare with ``depth'' \(\rightarrow\) \(\phi\)
	\item \textcolor{Green4}{see if \(\eta_k\) can be bounded in exact case - yes for the class of functions mentioned in the paper}
	\item \textcolor{orange}{find counterexample, that \(\eta\) can't be bounded in inexact case??? - main argument: have to assure, that convexified objective function is ``convex on all bundle points \(x^j\)'' from a certain \(\eta^k\) on }
	%\item \textcolor{red}{find out about \(\eta > \rho\)-condition? Or} \textcolor{Green4}{unnecessary}
	\item \textcolor{Green4}{compare nonconvex exact \(\leftrightarrow\) inexact convergence results \\
	only look at exact paper again if better results!}
	\begin{itemize}
		\item \textcolor{Green4}{check if correct: inexact more general because choice of \(t_k\) more freely??? in exact \(\mu_k\) only changed when restart \\
		update strategy not important for convergence; maybe for convergence speed?}
		\item check if update strategy important for convergence speed? - yes see napsu ... ``Comparison ...''
	\end{itemize}
	\item \textcolor{Green4}{check if (ii) in Theorem 6 in paper can be assured by choice of \(t_k\) in algorithm (think yes)}
	\item \textcolor{Green4}{compare to convergence results of other papers \\
	check if better results can be carried over - results all the same; check for prerequisites on functions}
	\item \textcolor{Green4}{check if other papers have better prerequisites \\
	check if results can be carried over \\
	all need locally Lipschitz and either a compact? subset or bounded lower level sets \\
	Better in other papers?: neither \(\{j \in J_k | \alpha_j^k >0\}\) nor \(\{\eta_k\}\) need to be bounded?? \\
	any prerequisites on \(t_k\) in other papers???}
	\item \textcolor{Green4}{most other solutions for nonconvex functions: based on dual idea; remark on dual view on bundle method?}
	\item \textcolor{Green4}{write that down nicely}
	\item read depth paper
	\begin{itemize}
		\item see what kind my algorithm is
		\item check Lemma 5 (iv) stronger?
	\end{itemize}
	\item remark on that inexactness makes objective nonconvex \(\rightarrow\) relate to nonconvexity paper?
	\item check again uniformly bounded \(J_k\) \\
		 aggregated objects (3.4), 5.2, 7.1 in ``depth''
	\item \textcolor{Green4}{generalized gradients may only be shifted not tilted? \\
		is this realistic assumption? \(\rightarrow\) all ok, tilted and shifted!}
		\item \textcolor{Green4}{ask Simon to \(\varepsilon\)-subdifferentials}
	\item Algorithm:
	\begin{itemize}
		\item print something useful in every iteration
		\item \(\eta\) should not get too big
	\end{itemize}
\end{itemize}

\textcolor{red}{If (newer) papers needed: look at citations of the ones I have.} \\

\textcolor{red}{(sub-)Level sets of continuous functions are closed (image of closed set closed under continuous function); should also hold for lower semicontinuous functions in metrizable(???) space}

\section{How is inexact information dealt with?}

\section{Extension with Limited Memory approach}

Variable metric methods are also known as quasi-Newton methods. \\

PhD Thesis (p. 33) inheritance of positive definiteness of matrix update formulas \\
p. 34: number of correction pairs usually  \(3\leq m \leq 30\)



\subsubsection{Thoughts about line search}
\begin{itemize}
	\item after what is written in ``nonconv, inex''-paper: Line search not provable to be finite if inexact information
	\item is line search standard in all variable metric methods?
	\begin{itemize}
		\item looks like it
		\item there do exist versions without line search \(\rightarrow\) other update???
		\item does is work without line search??? -think yes
	\end{itemize}
	\item does prox-parameter \(t_k\) have some relation to linesearch/stepsize?
\end{itemize}

\textbf{Algorithm 1 in \cite{Haarala2007}}
\vspace{1em}

\hrule  \vspace{0.4ex} \hrule
\vspace{1ex}
\textbf{Variable metric limited memory bundle method with inexact information}
\vspace{1ex}
\hrule
\vspace{1ex}
Select parameters \( m \in (0,1),~\gamma > 0,~ K>0,~ \rho \in (0,~ \frac{1}{2}) \) and a stopping tolerance \( \mathtt{tol} \geq 0\). \\
Set the initial metric matrix \(D_1=\mathbb{I}_{n\times n}\). Choose a starting point \(x^1 \in \R^n\) and compute \(f_1\) and \(s^1 = g^1 = S^1\). Set the initial serious iterate \(\hat{x}^1 := x^1\), \(\hat{f}_1 = f_1\), and \(\hat{s}^1 = s^1\), and \(c_1 = C_1 = 0\). \\
\textcolor{orange}{Set the correction indicator and the correction indicator for consecutive null steps \(i_C = 0\).}

For \(k = 1,2,3,  \dotsc \)   

\begin{enumerate}
	\item Compute \[d^k = -D_k S^k\]
	by using a limited memory BFGS update if the step before was a serious step and by using a limited memory SR1 update otherwise. \\
	For \(k=1\): \(d^1 = -S^1\).
	\item If \(-{S^k}^{\top}d^k < \rho {S^k}^{\top}S^k\) or \(i_{CN} = 1\) then set
	\begin{equation}
		d^k = d^k - \rho S^k,
	\end{equation}
	and \(i_C=1\). If the previous step was a null step set also \(i_{CN} = 1\).\\
	Otherwise set \(i_C = 0\).
	\item Set 
		\begin{align}
			\delta^w_k &= -{S^k}^{\top}d^k+2C_k \quad \text{ and} \\
			\delta^q_k &= \frac{1}{2} {S^k}^{\top}S^k + C_k.
		\end{align}
	If  \(\delta^w_k < \mathtt{tol}\) and \(\delta^q_k < \mathtt{tol}\) stop with \(\hat{x}_k\) as the final solution.
	\item Set 
		\begin{align}
		  \theta_k &= \min\{1, K/\|d^k\|\} \\
			x^{k+1} &= \hat{x}^k+\theta_kd^k.
		\end{align}
	Calculate the inexact function value \(f_{k+1}\) and an inexact subgradient \(g^{k+1}\).
\item If 
	\[f^{k+1} \leq \hat{f}^k - m\textcolor{red}{\delta^w_k} \quad \rightarrow \text{ serious step} \]
	\textcolor{red}{Which \(\delta\)???} \\
	Set \(\hat{x}^{k+1} = x^{k+1}, \hat{f}_{k+1} = f_{k+1}\), \(s^{k+1} = S^{k+1} = g^{k+1}\).\\
	Otherwise \(\rightarrow\) nullstep \\
	Compute the convexification parameter 
\[ \eta_{k+1} = \max\{e_{k+1}, 0\} + \gamma \]
with \(e_{k+1} = \hat{f}_{k+1}-f_{k+1}+\langle g^{k+1},d^k\rangle\). \\
	Set 
	\begin{align}
		\hat{x}^{k+1} &= \hat{x}^k \\
		\hat{f}^{k+1} &= \hat{f}^{k}\\
		s^{k+1} &= g^{k+1}+\eta_{k+1}*(x^{k+1}-\hat{x}^{k+1}) \\
		c_{k+1} &= \max\{e_{k+1},0\} + \frac{\gamma}{2}.
	\end{align}
	
	Compute the new correction pair \(u_1^k = \theta_kd^k\) and \(u_2^k = s^{k+1}-\hat{s}^{k+1}\).
	\item If this step was a serious step, go to 1. \\
	In case of a null step determine multipliers \(\alpha_i ^k\geq 0,~ i = \{1,2,3\},~ \sum_i{\alpha^k_i} =1\) that minimize the function
\begin{eqnarray*}
	\phi(\alpha_1, \alpha_2,\alpha_3) &=& (\alpha_1\hat{s}^{k+1}+\alpha_2s^{k+1}+\alpha_3S^{k})D_k(\alpha_1\hat{s}^{k+1}+\alpha_2s^{k+1}+\alpha_3S^{k}) \\
	&&+ 2(\alpha_2c_{k+1}+\alpha_3C_{k})
\end{eqnarray*}
where \(D_k\) is calculated by the same updating formula as in step 1 and \(D_k = D_k + \rho \mathbb{I} \) if \(i_C=1\). \\
Compute the aggregate subgradient and error

\begin{align}
	S^{k+1} &= \alpha_1\hat{s}^{k+1}+\alpha_2s^{k+1}+\alpha_3S^{k} \\
	C_{k+1} & = (\alpha_2c_{k+1}+\alpha_3C_{k})
\end{align}
and go back to step 1.
\end{enumerate}
\vspace{1ex}
\hrule

\vspace{1.5em}

\subsection{Convergence}



%
%
%
%\addcontentsline{toc}{section}{Bibliography}
%\bibliography{Bibliography}
	%\bibliographystyle{plain}
%
%
%\end{document}
\section{Variable Metric Bundle Algorithm}
\label{variable_metric}

...
We are now interested to include also curvature information into the method as soon as the objective function provides it.... \\

The proximal stabilization term mimics a curvature term pretty roughly; idea of variable metric bundle methods: allow not only Euclidean norm but a general norm \(\|x\cdot\|^2_{A} = \langle x,Ax\rangle\) with a symmetric and positive definite matrix \(A\).

This is also used by Noll et al. in \cite{Noll2012} and Noll in \cite{Noll2013} for nonconvex functions with exact and inexact information respectively.
In the original papers the motivation is to approximate the original objective function by a quadratic model

\begin{equation}
	\Phi(x,\hat{x}) = \phi(x,\hat{x})+\frac{1}{2}\langle x-\hat{x},Q(\hat{x})(x-\hat{x})\rangle
\end{equation}

with the convex and possibly nonsmooth first order model \(\phi(\cdot,\hat{x})\) and the quadratic but not necessarily convex second order part \(\frac{1}{2}\langle \cdot-\hat{x},Q(\hat{x})(\cdot-\hat{x})\rangle\).
This model is then again approximated to solve the bundle subproblem.

Here we take the same idea but interpret it more in the sense of a stabilization. The objective function is still first convexified and then this convexification approximated by a cutting plane model. The matrix \(Q(\hat{x})\) is added into the stabilization so that the \(k\)'th bundle subproblem is

\begin{equation}
	\min_{\hat{x}^k+d \in X} M_k(\hat{x}^k+d) + \frac{1}{2}\langle d,\left(Q+\frac{1}{t_k}\mathbb{I} \right) d \rangle
\end{equation}

...???


In \cite{Noll2012} and \cite{Noll2013} it is not specified how the matrix \(Q(\hat{x}\) is to be chosen.
For convergence it is necessary that the eigenvalues of \(Q(\hat{x}\) are bounded. We adopt here the notation in \cite{Noll2013} and say

\begin{equation}
	\exists q > 0 \text{ such that } -q\mathbb{I} \prec Q(\hat{x}) \prec q\mathbb{I}.
\end{equation}

Here \(A \prec B \) means that the matrix \(B-A\) is positive definite.

Additionally the matrix \(Q(\hat{x}+\frac{1}{t_k}\mathbb{I}\) has to be positive definite.

There are no other conditions put on \(Q(\hat{x}\).

It is not stated explicitly how to form such a matrix \(Q(\hat{x})\). Here we take an idea developed by Vl\v{c}ek  and Luk\v{s}an for a variable metric bundle method in \cite{Vlcek2001}.

The matrix \(Q(\hat{x}\) is formed by a BFGS-Update from the \textcolor{red}{last known - which? how many? all from bundle??} subgradients.



\begin{itemize}
	\item \(Q\) only updated in serious steps - why?
	\item 
\end{itemize}


\textcolor{red}{possible to do also SR1-update in null steps???}

\subsubsection{how to get Q and ensure all conditions}

\subsubsection{the new stopping criterion}

There are some minor changes that have to be made compared to the algorithm proposed by Hare et al. the biggest being the stopping condition.

In the same way as for (???) from the optimality condition

\begin{align}
	& 0 \in \partial M_k(x^{k+1})+\partial\mathtt{i}_{D}(x^{k+1})+\left(Q+\frac{1}{t_k}\mathbb{I}\right)d^k \\
	\label{Noll_opt_cond}
\end{align}

follows that 
\begin{equation}
	S^k+\nu^k = -\left(Q+\frac{1}{t_k}\mathbb{I}\right)d^k.
\end{equation}

From this the model decrease (???) can be recovered using (???), (??) and (\ref{Noll_opt_cond}):

\begin{align}
	\delta_k  &= \hat{f}_k - M_k(x^{k+1}) - (\nu^k)^{\top}d^k\\
	&= \hat{f}_k - A_k(x^{k+1}) - (\nu^k)^{\top}d^k\\
	%&= C_k - (S^k)^{\top}d^k - (\nu^k)^{\top}d^k \\
	&= C_k - (S^k+\nu^k)^{\top}d^k \\
	&= C_k + (d^k)^{\top}\left(Q+\frac{1}{t_k}\mathbb{I}\right)d^k
	\label{Noll_delta}
\end{align}







\subsection{Keywords}

%important in Noll for me: optimize model + \(d^{\top}(Q+\frac{1}{t_k}\mathbb{I})d\)
%-> some kind of second order information

%important: \(Q+\frac{1}{t_k}\mathbb{I}\) must have all eigenvalues \(\geq 0\).

%idea to get \(Q\): BFGS like in Fin-papers; theory

%!!! check stopping criterion
%connection between \(d^k\) and \(G^k/S^k\) now:
%Optimality condition: 


\subsection{important assumptions}
eigenvalues of \(Q\) are bounded \(\to\) possible by manipulating BFGS update

\begin{align}
	\text{if} &\quad norm\left(\frac{y^k{y^k}^{\top}}{{y^k}^{\top}d^k}\right) > 10^{??} \\
	&\quad \text{set } \frac{y^k{y^k}^{\top}}{threshold} \\
	&\quad threshold = norm\left(y^k{y^k}^{\top}\right)/10^{??} \\
	\text{end} &
\end{align}
same procedure for next term; all \(<1/3C\) for some overall threshold \(C\) 

\(Q+\frac{1}{t_k}\mathbb{I}\) such that \(\succ \xi \mathbb{I}\) for some fixed \(\xi > 0\). \\

\begin{equation}
	\min_{\hat{x}+d \in D} M^k(\hat{x}^k+d^k)+d^{\top}\frac{1}{2}\left(Q+\frac{1}{t_k}\mathbb{I}\right)d	
\label{Q_subprob}
\end{equation}
 

\subsection{Algorithm}
\vspace{1em}

\hrule  \vspace{0.4ex} \hrule
\vspace{1ex}
\textbf{Nonconvex variable metric bundle method with inexact information}
\vspace{1ex}
\hrule
\vspace{1ex}
Select parameters \( m \in (0,1), \gamma > 0 \) and a stopping tolerance \( \mathtt{tol} \geq 0\). \\
Choose a starting point \(x^1 \in \R^n\) and compute \(f_1\) and \(g^1\). Set the initial metric matrix \(Q = \mathbb{I}\), the initial index set \(J_1:=\{1\}\) and the initial prox-center to \(\hat{x}^1 := x^1\), \(\hat{f}_1 = f_1\) and select \(t_1 > 0\).

For \(k = 1,2,3,  \dotsc \)   

\begin{enumerate}
	\item Calculate \[d^k = \arg \min_{d \in \R^n} \left\{ M_k(\hat{x}^k+d)+\mathbb{I}_X(\hat{x}^k+d)+\frac{1}{2}d^{\top}\left(Q+\frac{1}{t_k}\mathbb{I}\right)d\right\}.\]
	\item Set %\textcolor{red}{\(\to\) other stopping condition!!!
		\begin{align*} 
		  G^k &= \sum_{j \in J_k}{\alpha_j^k s_j^k} \\ %, \quad	\nu^k = -\frac{1}{t_k}d^k-G^k???????????\\
			C_k &= \sum_{j \in J_k}{\alpha_j^k c_j^k} \\
	    \delta_k &=  C_k + (d^k)^{\top}\left(Q+\frac{1}{t_k}\mathbb{I}\right)d^k
		\end{align*}
		If \(\delta_k \leq \mathtt{tol} \rightarrow \) STOP.
	\item Set \( x^{k+1} = \hat{x}^k + d^k \).
	\item Compute \(f^{k+1}, g^{k+1}\) \\
	If 
	\[f^{k+1} \leq \hat{f}^k - m\delta_k \quad \rightarrow \text{ serious step} \]
	Set \(\hat{x}^{k+1} = x^{k+1}, \hat{f}^{k+1} = f^{k+1}\) and select \(t_{k+1} > 0\). \\
	Otherwise \(\rightarrow\) nullstep \\
	Set \(\hat{x}^{k+1} = \hat{x}^k, \hat{f}^{k+1}=f^{k+1}\) and choose \(0 < t_{k+1} \leq t_k\). 	
	\item Select new bundle index set \(J_{k+1}\), keeping all active elements. Calculate 
	\[ \eta_k \geq \max{\left\{\max_{j \in J_{k+1}, x^j \neq \hat{x}^{k+1}}{\frac{-2e_j^k}{|x^j - \hat{x}^{k+1}|^2}, 0}\right\}}+\gamma  \]
	and update the model \(M^k\)
\end{enumerate}
\vspace{1ex}
\hrule

\vspace{1.5em}

Lemma 5 in \cite{Hare2016} stays the same; no \(Q\) involved \\

\begin{theorem}(c.f.\cite[Theorem 6]{Hare2016}\\
	\(\to\) take only part with \(\liminf_{k\to\infty}t_k > 0 \) because other one not used in null steps and algorithm this way. \\
	Let the algorithm generate and infinite number of serious steps. Then \(\delta_k \to 0\) as \(k \to \infty\). \\
	Let the sequence \(\{\eta_k\}\) be bounded. If \(\liminf_{k \to \infty} t_k > 0\) then as \(k \to \infty\) we have \(C_k \to 0\), and for ever accumulation point \(\bar{x}\) of \(\{\hat{x}^k\}\) there exists \(\bar{S}\) such that \(S^k \to \bar{S}\) and \(S^k + \nu^k \to 0\). \\
	In particular if the cardinality of \({j \in J^k|\alpha_j^k > 0}\) is uniformly bounded in \(k\) then the conclusions of Lemma 5 in \cite{Hare2016} hold.
\end{theorem}

The proof is very similar to the one stated in \cite{Hare2016} but minor changes have to be made due to the different formulation of the nominal decrease \(\delta_k\).

\begin{proof}
	At each serious step \(k\) holds
	
	\begin{equation}
		\hat{f}_{k+1} \leq \hat{f}_k - m\delta_k
	\label{nonincreasing}
	\end{equation}
	
	where \(m, ~\delta_k > 0\). From this follows that the sequence \(\{\hat{f}_k\}\) is nonincreasing.
	Since \(\{\hat{x}^k\} \subset D\) the sequence is by the fact that \(f\) is finite %\textcolor{red}{which assumption says \(f\) bounded below??? \(\to\) \(f\) subdifferentially regular \(\Rightarrow\) \(f\) finite}
	and \(|\sigma_k| < \bar{\sigma}\) the sequence \(\{f(\hat{x}^k)+\sigma_k\} = \{\hat{f}_k\}\) is bounded below. Together with the fact that \(\{\hat{f}_k\}\) is nonincreasing one can conclude that it converges. \\
	Using (\ref{nonincreasing}), one obtains
	
	\begin{equation}
		0 \leq m \sum_{k = 1}^l \delta_k \leq \sum_{k = 1}^l \left(\hat{f}_k-\hat{f}_{k+1}\right),
	\end{equation}
	
	so letting \(l \to \infty\), 
	
	\begin{equation}
		0 \leq m\sum_{k=1}^{\infty} \delta_k \leq \hat{f}_1 - \underbrace{\lim_{k \to \infty} \hat{f}_k}_{\neq \pm \infty}.
	\end{equation}
	As a result,
	\begin{equation}
		\sum_{k = 1}^{\infty} \delta_k = \sum_{k=1}^{\infty}\left(C^k+(d^k)^{\top}\left(Q+\frac{1}{t_k}\mathbb{I}\right)d^k\right) < \infty
	\end{equation}
	
	Hence, \(\delta_k \to 0\) as \(k \to \infty\). As all quantities above are nonnegative due to positive (semi-)definiteness of \(Q+\frac{1}{t_k}\mathbb{I}\), it also holds that
	
	\begin{equation}
		C_k \to 0 \quad \text{and} \quad (d^k)^{\top}\left(Q+\frac{1}{t_k}\mathbb{I}\right)d^k \to 0.
	\end{equation}
	
	For any accumulation point \(\bar{x}\) of the sequence \(\{\hat{x}^k\}\) the corresponding subsequence \(d^k \to 0\) for \(k \in K \subset \{1,2,...\} \). As \(\liminf_{k \to \infty} t_k > 0\) and the eigenvalues of \(Q\) are bounded the whole expression 
	
	\begin{equation}
	 S^k + \nu^k = \left(Q+\frac{1}{t_k}I \right)d^k  \to 0 \quad \text{for} \quad k \in K.
	\end{equation}
	
	And from local Lipschitz continuity of \(f\) follows then that \(S^k \to \bar{S}\) for \(k \in K\).
	
\end{proof}

\begin{remark}
If one assumes that the set \(\Omega = \{x \in \R^n |f(x) \leq f(x^1) +2\bar{\sigma}\}\) is bounded, it is not necessary to use the constraint set \(D\). \\
Because all \(\{\hat{x}^k\} \subset \Omega\) one can deduce the boundedness of the sequence.
\end{remark}

For the case of infinitely many null steps one show:

\begin{theorem} \cite{Hare2016}
	Let a finite number of serious iterates be followed by infinite null steps. Let the sequence \(\{\eta_k\}\) be bounded and \(\liminf k \to \infty > 0\). \\
	Then \(\{x^k\} \to \hat{x}\), \(\delta_k \to 0\), \(C_k \to 0\), \(S^k + \nu^k \to 0\) and there exist \(K\subset \{1,2,...\}\) and \(\bar{S}\) such that \(S^k \to \bar{S^k}\) as \(K \ni k \to \infty\). \\
	In particular if the cardinality of \({j \in J^k|\alpha_j^k > 0}\) is uniformly bounded in \(k\) then the conclusions of Lemma 5 in \cite{Hare2016} hold for \(\bar{x} = \hat{x}\). 
\end{theorem}

\begin{proof}
	Let \(k\) be large enough such that \(k \geq \bar{k}\) and \(\hat{x}^k = \hat{x}\) and \(\hat{f}_k=\hat{f}\) are fixed.
	Define the optimal value of the subproblem (\ref{Q_subprob}) by 
	
	\begin{equation}
		\Psi_k := M_k(x^{k+1})+\left(d^k\right)^{\top}\frac{1}{2}\left(Q+\frac{1}{t_k}\mathbb{I}\right)d^k.
		\label{Psi}
	\end{equation}
	
	It is first shown that the sequence \(\{\Psi_k\}\) is bounded above.
	From the definition

\begin{equation}	
	A_k(\hat{x}+d) := M_k(x^{k+1})+\langle S^k,d-d^k\rangle
	\label{Noll_agg_lin}
	of the aggregate linearization follows
	
	\begin{equation}
		A_k(\hat{x}) = M_k(x^{k+1})-\langle S^k,d^k \rangle.
	\end{equation}
	
	Using \(S^k+\nu^k = -\left(Q+\frac{1}{t_k}\mathbb{I}\right)d^k\) and the subgradient inequality for \(\nu^k \in \partial \mathtt{i}_{D}\) one obtains
	
	\begin{align*}
		\Psi^k+\frac{1}{2}\left(d^{k}\right)^{\top}\left(Q+\frac{1}{t_k}\mathbb{I}\right)d^k &= A_k(\hat{x})+\langle S^k,d^k\rangle + \left(d^{k}\right)^{\top}\left(Q+\frac{1}{t_k}\mathbb{I}\right)d^k \\
		&= A_k(\hat{x})-\langle \nu^k,k \rangle \\
		&\leq A(\hat{x}) \\
		&\leq M_k(\hat{x}) \\
		& = \hat{f}
	\end{align*}
	
	\textcolor{red}{where the equations and inequalities follow from???}\\
	By boundedness of \(d^k\) and \(Q+\frac{1}{t_k}\mathbb{I}\) this yields that \(\Psi_k \leq \hat{f}\), so the sequence \(\{\Psi_k\}\) is bounded above.
	In the next step is shown that \(\{\Psi_k\}\) is increasing.
	
	\begin{align}
		\Psi_{k+1} &= M_k(x^{k+2})+\frac{1}{2}\left(d^{k+1}\right)^{\top}\left(Q+\frac{1}{t_k}\mathbb{I} \right)d^{k+1} \\
		&\geq A_k(x^{k+2})+\frac{1}{2}\left(d^{k+1}\right)^{\top}\left(Q+\frac{1}{t_k}\mathbb{I}\right)d^{k+1} \\
		&= M_k(x^{k+1})+\langle S^k,x^{k+2}-x^{k+1} \rangle +\frac{1}{2}\left(d^{k+1}\right)^{\top}\left(Q+\frac{1}{t_k}\mathbb{I}\right)d^{k+1} \\
		&= \Psi_k - \frac{1}{2}\left(d^{k}\right)^{\top}\left(Q+\frac{1}{t_k}\mathbb{I}\right)d^{k} + \frac{1}{2}\left(d^{k+1}\right)^{\top}\left(Q+\frac{1}{t_k}\mathbb{I}\right)d^{k+1} \\
		& \qquad -\left(d^{k}\right)^{\top}\left(Q+\frac{1}{t_k}\mathbb{I}\right)\left(d^{k+1} - d^{k}\right) - \langle \nu^k, x^{k+2}-x^{k+1}\rangle\\
		&\geq \Psi_k + \frac{1}{2}\left(d^{k+1}-d^{k}\right)^{\top}\left(Q+\frac{1}{t_k}\mathbb{I}\right)\left(d^{k+1} - d^{k}\right)
	\end{align}
	
\textcolor{red}{say where (in-)equalities come from} \\

As \(Q\) is fixed in null steps and \(\liminf_{k \to \infty} t_k > 0\) \(\{\Psi_k\}\) is increasing. The sequence is therefore convergent.
Consequently, taking into account that \(1/t_k \geq 1/t_{\bar{k}}\), it follows

\begin{equation}
	\|d^{k+1}-d^k\| \to 0, \quad k \to \infty.
\label{d_to_0}
\end{equation}

By the \textcolor{red}{definitions and characterizations that have to be specified} one has

\begin{align}
	\hat{f} &= \delta_k+M_k(\hat{x})-C_k-\left(d^{k}\right)^{\top}\left(Q+\frac{1}{t_k}\mathbb{I}\right)\left(d^{k}\right) \\
	&= \delta_k + M_k(x^{k+1})-\langle S^k,d^k \rangle-\left(d^{k}\right)^{\top}\left(Q+\frac{1}{t_k}\mathbb{I}\right)\left(d^{k}\right)  \\
	&= \delta_k \\
	&\geq \delta_k+M_k(\hat{x}+d^k)
\end{align}

Where the last inequality is given by \(\nu^k \in \partial\mathtt{i}_{D}(x^{k+1})\). Therefore 

\begin{equation}
	\delta^{k+1} \leq \hat{f}-M_{k+1}(\hat{x}+d^{k+1}).
\end{equation}

By the first inequality in assumption \textcolor{red}{define assumption} on the model, written for \(d=d^{k+1}\),

\begin{equation}
	-\hat{f}_{k+1}+c^{k+1}_{k+1}-\left\langle s^{k+1}_{k+1},d^{k+1}\right\rangle \geq -M_{k+1}(\hat{x}+d^{k+1}).
\end{equation}
As \(\hat{f}_{k+1}=\hat{f}\), adding condition \textcolor{red}{???} to the inequality above, one obtains that

\begin{equation}
	m\delta_k+\left\langle s^{k+1}_{k+1},d^k-d^{k+1}\right\rangle \geq \hat{f}-M_{k+1}(\hat{x}+d^{k+1}).
\end{equation}

Combining this relation with \textcolor{red}{???} yields

\begin{equation}
	0 \leq \delta_{k+1} \leq m\delta_k + \left\langle s^{k+1}_{k+1},d^k-d^{k+1}\right\rangle.
	\label{delta_ineq}
\end{equation}

Since \(m \in (0,1)\) and \(\left\langle s^{k+1}_{k+1},d^k-d^{k+1}\right\rangle \to 0\) as \(k \to \infty\) due to (\ref{d_to_0}) and the boundedness of \(\{\eta_k\}\) using \cite[Lemma 3, p.45]{Polyak1987} it follows from  (\ref{delta_ineq}) that 

\begin{equation}
	\lim_{k \to \infty} \delta_k = 0.
\end{equation}

From the formulation \(\delta_k = C_k + \left(d^k\right)^{\top}\left(Q+\frac{1}{t_k}\mathbb{I}\right)d^k\) follows that \(C_k \to 0\) as \(k \to \infty\). As \textcolor{red}{\(Q+\frac{1}{t_k}\mathbb{I} \succ \xi\mathbb{I}\)} it follows that 

\begin{equation}
	\xi \left(d^k\right)^{\top}d^k \leq \left(d^k\right)^{\top}\left(Q+\frac{1}{t_k}\mathbb{I}\right)d^k \to 0
\end{equation}

\end{proof}

\begin{remark}
	To the boundedness of \(Q\) and the resulting inequalities:
	\[ A \succ B \Leftrightarrow A - B \succ 0\Leftrightarrow A-B \text{ is positive definite}\]
	Respectively for \(A \prec B\) \\
  For a real symmetric matrix \(A\) and a vector \(d\in \R^n\) the following result holds: 
	\[ A \prec \xi \mathbb{I} \Rightarrow Ad < \xi d \]
	\textbf{Proof:} \(A\) real and symmetric \(\Rightarrow\) it is orthogonally diagonalizeable:
	\begin{equation}
	\begin{split}
		\exists \lambda_i \in \R \text{ eigenvalues} \quad v^i \in \R^n \text{ eignvektors} \\
		Av^i= \lambda_iv^i \\
		\text{and} \quad \exists \alpha_i \in \R: \quad d = \sum_{i} {\alpha_i v^i} \\
		\Rightarrow Ad = A \sum_i{\alpha_i v^i} = \sum_i \alpha_i Av^i = \sum_i \alpha_i \lambda_i v^i\\
		A \prec \xi \mathbb{I}  \Leftrightarrow \max_i \lambda_i < \xi \\
		\Rightarrow Ad < \xi \sum_i \alpha_i v^i = \xi d
  \end{split}
	\end{equation}
\end{remark}












%\subsection{Subproblem Variable Metric}
%
%For comparision: Subproblem proximal bundle
%
%\begin{align}
	%&\min_{d \in \R^n, \xi \in \R} \xi + \frac{1}{2t_k}\|d\|^2 = \xi + \frac{1}{2}d^{\top}\left(\frac{1}{t_k}\mathbf{I}\right)d \\
	%&\text{s.t.} \quad f(\hat{x}^k)+{g^j}^{\top}d - e^k_j - \xi \leq 0, \quad j \in J_k
%\end{align}
%
%Subproblem variable metric:
%\begin{align}
	%&\min_{d \in \R^n, \xi in \R} \xi + \frac{1}{2}d^{\top}D_kd \\
	%&\text{s.t.} \quad f(\hat{x}^k)+{g^j}^{\top}d - e^k_j - \xi \leq 0, \quad j \in J_k
%\end{align}
%These are \(\R^{n+1}\) dimensional quadratic optimization problems.\\
%
%\textcolor{red}{Find out if \(D_k\) is diagonal matrix! Think not.} \\
%
%Approaches not so different. Instead of just scaling the identity \(\rightarrow\) induce ``curvature information'' via past subgradients. \\
%
%Dual proximal subproblem:
%\begin{align}
	%&\min_{\alpha \in \R^{|J_k|}} \frac{1}{2} \left(\sum_{j \in J_k}{\alpha_jg^j}\right)^{\top} t_k\mathbf{I} \left(\sum_{j \in J_k}{\alpha_jg^j}\right) + \sum_{j \in J_k}{\alpha_j e_j^k} \\
		%&\text{s.t.} \quad \sum_{j \in J_k}{\alpha_j} = 1 \text{ and } \alpha_j \geq 0 ~ j \in J_k
%\end{align}
%
%
%Dual variable metric subproblem:
%\begin{align}
	%&\min_{\alpha \in \R^{|J_k|}} \frac{1}{2} \left(\sum_{j \in J_k}{\alpha_jg^j}\right)^{\top} D^{-1}_k \left(\sum_{j \in J_k}{\alpha_jg^j}\right) + \sum_{j \in J_k}{\alpha_j e_j^k} \\
		%&\text{s.t.} \quad \sum_{j \in J_k}{\alpha_j} = 1 \text{ and } \alpha_j \geq 0 ~ j \in J_k
%\end{align}
%
%These are \(\R^{|J_k|}\) dimensional quadratic optimization problems. \\
%
%\textcolor{red}{check linear independent \(g^j\)'s.}

%\section{Application to Model Selection for Primal SVM}

\textcolor{red}{Skalarprodukt anpassen, Vektoren nicht fett oder neue definition, notation, \(\lambda \in \Lambda\) einfugen}

\subsection{Introduction}

In this chapter the nonconvex inexact bundle algorithm is applied to the problem of model selection for \emph{support vector machines} (SVM) solving classification tasks.
It relies on a bilevel formulation proposed by Kunapuli \cite{Kunapuli2008} and Moore et al. \cite{Moore2011}.

A natural application for the inexact bundle algorithm is an optimization problem where the objective function value can only be computed iteratively. This is for example the case in bilevel optimization.
%If this objective function is iteratively optimized, the function value i never exact. But by the stopping condition it is (often???) possible to have the needed error bound on the function value (what about the gradient???). (Comment on asymptotical exactness?)

%A general bilevel problem can be formulated as 
%Bilevel optimization deals with a class of constrained optimization problems where one or more of the optimization variables are again constrained by optimization problems.
A general bilevel program can be formulated as \cite{Kunapuli2008}
\begin{equation}
	\begin{array}{cll}
	\displaystyle\min_{x \in X, y} & F(x,y) & \text{upper level} \vspace{0.5ex}\\
	\text{s.t.} & G(x,y) \leq 0 & \vspace{1em}\\
	& y \in \begin{Bmatrix} \displaystyle\argmax_{y\in Y} & f(x,y) \vspace{0.5ex}\\
	                        \text{s.t.} & g(x,y) \leq 0 
													\end{Bmatrix}. & \text{lower level}
	\end{array}
\end{equation}

It consists of an \emph{upper} or \emph{outer level} which is the overall function to be optimized. Contrary to usual constrained optimization problems which are constrained by explicitly given equalities and inequalities a  bilevel program is additionally constrained to a second optimization problem, the \emph{lower} or \emph{inner level} problem.

Solving bilevel problems can be divided roughly in two classes: implicit and explicit solution methods. \\
In the explicit methods the lower level problem is usually rewritten by its KKT conditions and the upper and lower level are solved simultaneously. For the setting of model selection for support vector machines as it is used here, this method is described in detail in \cite{Kunapuli2008}.

The second approach is the implicit one. Here the lower level problem is solved directly in every iteration of the outer optimization algorithm and the solution is plugged into the upper level objective. \\
Obviously if the inner level problem is solved numerically, the solution cannot be exact. Additionally the \emph{solution map} \( S(x) = \{y \in \R^k | y\) solves the lower level problem\(\}\) is often nondifferentiable \cite{Outrata1998} and since elements of  the solution map are plugged into the outer level objective function in the implicit approach, the outer level function becomes nonsmooth itself. \\
This is why the inexact bundle algorithm seems a natural choice to tackle these bilevel problems. \\ 
Moore et al. use the implicit approach in \cite{Moore2011} for support vector regression. However they use a gradient decent method which is not guaranteed to stop at an optimal solution. \\
In \cite{Moore2010a} he also suggests the nonconvex exact bundle algorithm of Fuduli et al. \cite{Fuduli2004a} for solving the bilevel regression problem. This allows for nonsmooth inner problems and can theoretically solve some of the issues of the gradient descent method. It ignores however, that the objective function values can only be calculated approximately. A fact which is not addressed in Fuduli's algorithm.

%\subsection{Notation and Expressions}

%training set, validation /holdout set...

\subsection{Introduction to Support Vector Machines}
%In times of big data machine learning is a very active field of research. 
Support vector machines are linear learning machines that were developed in the 90's by Vapnik and co-workers. Soon they could outperform several other programs in this area \cite{Cristianini2000} and the subsequent interest in SVMs lead to a very versatile application of these machines \cite{Kunapuli2008}.

The case that is considered here is binary support vector classification using supervised learning. \\
In classification data from a possibly high dimensional vector space \(\tilde{X} \subseteq \R^n\), the \emph{feature} or \emph{input space} is divided into two classes. These lie in the \emph{output domain} \(\tilde{Y} = \{-1,1\}\). Elements from the feature space will mostly be called \emph{data points} here. They get \emph{labels} from the feature space. Labeled data points are called \emph{examples}. \\
The functional relation between the features and the class of an example is given by the usually unknown \emph{response} or \emph{target function} \(f(x)\). \\
Supervised learning is a kind of machine learning task where the machine is given examples of input data with associated labels, the so called \emph{training data} \((X,Y)\). Mathematically this can be modeled by assuming that the examples are  drawn identically and independently distributed (iid) from the fixed joint distribution \(P(x,y)\). This usually unknown distribution states the probability that an data point \(x\) has the label \(y\) \cite{Vapnik1999}. \\
The overall goal is then to optimize the generalization ability, meaning the ability to predict the output for unseen data correctly \cite{Cristianini2000}.
%The overall goal of machine learning is to optimize the generalization ability of the algorithm --> this means: predict output for unlabeled input as well as possible.
%task is to find a mapping that predicts the output given unlabeled input as good as possible \cite{Cristianini2000}. 

\subsubsection{Risk minimization}
The concept of SVM's was originally inspired by the statistical learning theory developed by Vapnik. For a throughout analysis see \cite{Vapnik1998}. \\
%from: Vapnik - Statistical Learning Theory overview \cite{Vapnik1998}
The idea of \emph{risk minimization} is to find from a fixed set or class of functions the one that is the best approximation to the response function. This is done by minimizing a loss function that compares the given labels of the examples to the response of the learning machine.
%goal: find best available (from set of functions that was chosen before) approximation to the actual (unknown) ``response function'' --> function that ``classifies'' the data
%Minimize loss between response of ``supervisor'' and response of learning machine

As the response function is not known only the expected value of the loss can be calculated. It is given by the \emph{risk functional} 

\begin{equation}
	R(\lambda) = \int{\mathcal{L}(y,f_{\lambda}(x))\text{d}P(x,y)}
\label{risk_func}
\end{equation}

Where \(\mathcal{L}: \R^2 \to \R\) is the loss function, \(f_{\lambda}: \R^{n}\cap\mathcal{F} \to \R, ~ \lambda \in \Lambda\) the response function found by the learning machine and \(P(x,y)\) the joint distribution the training data is drawn from. The goal is now to find a function \(f_{\bar{\lambda}}(x)\) in the chosen function space \(\mathcal{F}\) that minimizes this risk functional \cite{Vapnik1999}.

As the only given information is given by the training set inductive principles are used to work with the empirical risk, rather than with the risk functional.
The empirical risk only depends on the finite training set  and is given by 

%goal: find \(f_{\lambda_0}(x)\) that minimizes this functional
%\(P(x,y)\) is unknown
%Because only a data sample is known: Work with the empirical risk

\begin{equation}
	R_{emp}(\lambda) = \frac{1}{l} \sum_{i = 1}^{l}\mathcal{L}(y_i,f_{\lambda}(x^i)),
\label{emp_risk}
\end{equation}

where \(l\) is the number of data points.
%L: Loss function 
%Goal: approximate \(Q(z,\lambda_0)\) by \(Q(z,\lambda_l)\) that minimizes the empirical risk. --> is called ERM (empirical risk minimization induction principle)
The law of large numbers ensures that the empirical risk converges to the risk functional as the number of data points grows to infinity. This however does not guarantee that the function \(f_{\lambda,emp}\) that minimizes the empirical rist also converges towards the function \(f_{\bar{\lambda}}\) that minimizes the risk functional.
The theory of consistency provides necessary and sufficient conditions that solve this issue \cite{Vapnik1999}.

Vapnik introduced therefore the structural risk minimization induction principle (SRM). It ensures that the used set of functions has a structure that makes it strongly consistent \cite{Vapnik1999}. Additionally it takes the complexity of the function that is used to approximate the target function into account. ``The SRM principle actually suggests a tradeoff between the quality of the approximation and the complexity of the approximating function'' \cite[p. 994]{Vapnik1999}.
This reduces the risk of \emph{overfitting}, meaning to overly fit the function to the training data with the result of poor generalization \cite{Cristianini2000}.

%Therefore: structural risk Minimization induction Principle (SRM) --> is strongly consistent --> idea: minimize the empirical risk by considering the complexity of the model (given by CV-dimension)

Support vector machines fulfill all conditions of the SRM principle. Due to the kernel trick that allows for nonlinear classification tasks it is also very powerful. For more detailed information on this see \cite{Kunapuli2008, Vapnik1998} and references therein.

\subsubsection{Support Vector machines}

In the case of linear binary classification one searches for a an affine hyperplane \(\bm{w}\) shifted by \(b\) to separate the given data. The vector \(\bm{w}\) is called weight vector and \(b\) is the bias. \\
Let the data be linearly separable. The function deciding how the data is classified can then be written as

\[ f(x) = sign(\bm{w}^{\top}x-b). \]

Support vector machines aim at finding such a hyperplane that separates also unseen data optimally.

\textcolor{red}{???Picture of hyperplane}

One problem of this intuitive approach is that the representation of a hyperplane is not unique. If the plane described by \((\bm{w},b)\) separates the data there exist infinitely many hyperplanes \((t\bm{w},b), ~ t>0\) that separate the data in the same way. \\
To have a unique description of a separation hyperplane the \emph{canonical hyperplane for given data} \(x \in X\) is defined by 

\[ f(x) = \bm{w}^{\top} x - b  \quad \text{s.t.} \quad \min_i |\bm{w}^{\top}x^i-b| = 1\] 

This is always possible in the case where the data is  linearly separable and means that the inverse of the norm of the weight vector is equal to the distance of the closest point \(x \in X \) to  the hyperplane \cite{Kunapuli2008}.


%Kunapuli: representation of hyperplane not unique --> if \((\bm{w},b)\) separates data, infinitely many hyperplanes \((t\bm{w},tb) \forall t > 0\) also separate data --> need unique representation --> w.l.o.g. define \emph{canonical hyperplane} for given data \(x \in X\) as function 
%only valid for separable data

This gives rise to the following definition:
The \emph{margin} is the minimal Euclidean distance between a training example \(x^i\) and the separating hyperplane.
%one can show that margin invese proportional to norm of \(\bm{w}\).
A bigger margin means a lower complexity of the function \cite{Cristianini2000}. 

A \emph{maximal margin hyperplane} is the hyperplane that realizes the maximal possible margin for a given data set.

%Theorem 6.1
\begin{theorem}[{\cite[Theorem 6.1]{Cristianini2000}}]
Given a linearly separable training sample \(\Omega = ((x^i,y_i),...,(x^l,y_l))\) the hyperplane \((\bm{w},b)\) that solves the optimization problem

\[\|\bm{w}\|^2 \quad \text{subject to} \quad y_i(\bm{w}^{\top}x-b)\geq 1 \quad i = 1,...,l \]

realizes a maximal margin hyperplane  
\end{theorem}

Generally one cannot assume the data to be linearly separable. This is why in most applications a so called \emph{soft margin classifier} is used. It introduces the slack variables \(\xi_i\) that measure the distance of the misclassified points to the hyperplane:

%from this follows: soft margin classifyer --> not any more assumption that data linearly separable --> slack variables \(\xi_i\) measure the distance of misclassified points

%Def 2.6:
%\begin{definition} \cite{Cristianini2000}
Fix \(\gamma > 0\). A \emph{margin slack variable of the example} \((x^i,y_i)\) with respect to the hyperplane \((\bm{w},b)\) and target margin \(\gamma\) is 

\[\xi_i = \max(0, \gamma - y_i(\bm{w}^{\top}x+b))\] 

%\end{definition}  
If \(\xi_i > \gamma\) the point is misclassified. \\
One can also say that \(\|\xi\|\) measures the amount by which training set ``fails to have margin \(\gamma\)'' \cite{Cristianini2000}.

For support vector machines the target margin is set to \(\gamma = 1\).

%Book uses \(\xi^2\) -> then \(\xi>0\) not needed any more; only \(\xi\) also used -> 1-norm
This results finally in the following slightly different optimization problems for finding an optimal separating hyperplane \((\bm{w},b)\): 

\begin{align}
\begin{split}
	\min_{\bm{w},b,\xi} \quad & \frac{1}{2} \|\bm{w}\|^2_2+C\sum_{i=1}^l{\xi_i} \\
	\text{subject to} \quad & y_i(\bm{w}^{\top}x^i-b) \geq 1 - \xi_i \\
	& \xi_i \geq 0 \\
	& 	\forall i = 1, \dots , l
\end{split}
\label{SVM_1}
\end{align}

and

\begin{align}
\begin{split}
	\min_{\bm{w},b,\xi} \quad & \frac{1}{2} \|\bm{w}\|^2_2+C\sum_{i=1}^l{\xi_i^2} \vspace{0.5ex}\\
	\text{subject to} \quad & y_i(\bm{w}^{\top}x^i-b) \geq 1 - \xi_i \vspace{0.5ex}\\
	& 	\forall i = 1, \dots , l
\end{split}
\label{SVM_2}
\end{align}


The parameter \(C > 0\) gives a trade-off between the richness of the chosen set of functions \(f_{\alpha}\) to reduce the error on the training data and the danger of overfitting to have good generalization. It has to be chosen a priori \cite{Kunapuli2008}.


\subsection{Explanation Bilevel Approach and Inexact Bundle Method}
%Check if regression case also included?????

The hyper-parameter \(C\) in the objective function of the classification problem has to be set before hand. This step is part of the model selection process.
To set this parameter optimally different methods can be used.
A very intuitive and widely used approach is doing  and \emph{cross validation} (CV) with a grid search implementation.
%very costly, discrete parameter choice, not practicable in case of many parameter

To prevent overfitting and get a good parameter selection, especially in case of little data, commonly \(T\)-fold cross validation is used. \\
For this technique the training data is randomly partitioned into \(T\) subsets of equal size. One of these subsets is then left out for training and instead used afterwards to get an estimate of the generalization error.  \\
To use CV for model selection it has to be embedded into an optimization algorithm over the hyper-parameter space. 
Commonly this is done by discretizing the parameter space and for \(T\)-fold CV training \(T\) models at each grid point. The resulting models are then compared to find the best parameters in the grid.
Obviously for a growing number of hyper-parameters this is very costly. An additional drawback is that the parameters are only chosen from a finite set \cite{Kunapuli2008}.

\subsubsection{Reformulation as bilevel problem}

A more recent approach is the formulation as a bilevel problem used in \cite{Kunapuli2008, Moore2011}.
This makes it possible to optimize the hyper-parameters continuously.

%Formal description of the bilevel problem for t-fold cross validation
%Better model description --> see 2.2 (p.46) in thesis 

Let \(\Omega = {(x^1,y_1),...,(x^l,y_l)} \subseteq \R^{n+1}\) be a given data set of size \(l = |\Omega|\). The associated index set is denoted by \(\mathcal{N}\). For classification the labels \(y_i\) are \(\pm1\).
For \(T\)-fold cross validation let \(\bar{\Omega}_t\) and \(\Omega_t\) be the training set and the validation set within the \(t\)'th fold and \(\bar{\mathcal{N}}_t\) and \(\mathcal{N}_t\) the respective index sets.
Furthermore let \(f^t:\R^{n+1}\cap\mathcal{F} \to \R\) be the response function trained on the \(t\)'th fold and \(\lambda \in \Lambda\) the hyper-parameters to be optimized.
For a general machine learning problem with upper and lower loss function \(\mathcal{L}_{upp}\) and \(\mathcal{L}_{low}\) respectively the bilevel problem writes
 
%%% introduce \mathcal(F) already before as the function space

\begin{equation}
	\begin{array}{cll}
	\displaystyle\min_{\lambda, f^t} & \mathcal{L}_{upp}\left(\lambda,f^1|_{\Omega_1},...,f^T|_{\Omega_T}\right) & \text{upper level} \vspace{0.5ex}\\
	\text{s.t.} & \lambda \in \Lambda & \vspace{1em}\\
	& \text{for } t = 1,...,T: & \\
	& f^t \in \begin{Bmatrix} \displaystyle\argmin_{f\in \mathcal{F}} & \mathcal{L}_{low}(\lambda,f,(x^i,y_i)_{i = 1}^l\in \bar{\Omega}_t) \vspace{0.5ex}\\
	                        \text{s.t.} & g_{low}(\lambda,f) \leq 0 
													\end{Bmatrix}. & \text{lower level}
	\end{array}
\end{equation}

In the case of support vector classification the \(T\) inner problems are one of the classical SVM formulations (\ref{SVM_1}) or (\ref{SVM_2}) (but all \(T\) problems have the same formulation).
The problem can also be rewritten into a unconstrained form. This form will be helpful when using the inexact bundle algorithm for solving the bilevel problem.
For the \(t\)'th fold the resulting hyperplane is identified with the pair \((\bm{w}^t,b_t) \in \R^{n+1}\).
The inner level problem for the \(t\)'th fold can therefore be stated as

\begin{equation}
	(\bm{w}^t,b_t) \in \argmin_{\bm{w},b}\left\{ \frac{\lambda}{2} \|\bm{w}\|^2_2+\sum_{i\in \bar{\mathcal{N}}_t}{\max\left(1-y_i(\bm{w}^{\top}x^i-b),0\right)}\right\}
\label{lower_t_1}
\end{equation}

or

\begin{equation}
	(\bm{w}^t,b_t) \in \argmin_{\bm{w},b} \left\{\frac{\lambda}{2} \|\bm{w}\|^2_2+\sum_{i\in \bar{\mathcal{N}}_t}{\left(\max\left(1-y_i(\bm{w}^{\top}x^i-b),0\right)\right)^2}\right\}
\label{lower_t_2}
\end{equation}

Where the hyper-parameter \(\lambda = \frac{1}{C}\) was used  due to numerical stability \cite{Kunapuli2008}.

For the upper level objective function there are different choices possible.
Simply put the outer level objective should compare the different inner level solutions and pick the best one. An intuitive choice would therefore be to pick the misclassification loss, that count how many data points of the respective validation set \(\Omega_t\) were misclassified when taking function \(f^t\).

The misclassification loss can be written as

\begin{equation}
	\mathcal{L}_{mis} = \frac{1}{T}\sum_{t=1}^T\frac{1}{|\mathcal{N}_t|}\sum_{i \in \mathcal{N}_t}{\left[-y_i((\bm{w}^t)^{\top}x-b_t)\right]_{\star}}
\label{misclass_loss}
\end{equation}

where the step function \(()_{\star}\) is defined componentwise for a vector as
\begin{equation}
	(r_{\star})_i = \left\{\begin{array}{c} 1, \quad \text{if } r_i > 0, \\ 0, \quad \text{if } r_i \leq 0 \end{array} \right. .
\label{step_fun}
\end{equation}

The drawback of this simple loss function is, that it is not continuous and as such not suitable for subgradient based optimization. Therefore another loss function is used for the upper level problem - the \emph{hinge loss}. It is an upper bound on the misclassification loss and reads

\begin{equation}
		\mathcal{L}_{hinge} = \frac{1}{T}\sum_{t=1}^T\frac{1}{|\mathcal{N}_t|}\sum_{i \in \mathcal{N}_t}{\max\left(1-y_i((\bm{w}^t)^{\top}x-b_t),0\right)}.
\label{hinge_loss}
\end{equation}

Hence the two final resulting bilevel formulations for model selection in support vector are 

\begin{align}
\begin{split}
	\min_{\bm{W},\bm{b}} \quad &  \mathcal{L}_{hinge}(\bm{W},\bm{b}) = \frac{1}{T}\sum_{t=1}^T\frac{1}{|\mathcal{N}_t|}\sum_{i \in \mathcal{N}_t}{\max\left(1-y_i((\bm{w}^t)^{\top}x-b_t),0\right)}\\
	\text{subject to} \quad &  \lambda > 0 \\
	& \text{for } t = 1,...,T \\
	& (\bm{w}^t,b_t) \in \argmin_{\bm{w},b}\left\{ \frac{\lambda}{2} \|\bm{w}\|^2_2+\sum_{i\in \bar{\mathcal{N}}_t}{\max\left(1-y_i(\bm{w}^{\top}x^i-b),0\right)}\right\} \\
\end{split}
\label{model_1}
\end{align}

and 

\begin{align}
\begin{split}
	\min_{\bm{W},\bm{b}} \quad &  \mathcal{L}_{hinge}(\bm{W},\bm{b}) = \frac{1}{T}\sum_{t=1}^T\frac{1}{|\mathcal{N}_t|}\sum_{i \in \mathcal{N}_t}{\max\left(1-y_i((\bm{w}^t)^{\top}x-b_t),0\right)}\\
	\text{subject to} \quad & \lambda > 0 \\
	& \text{for } t = 1,...,T \\
	& (\bm{w}^t,b_t) \in \argmin_{\bm{w},b}\left\{ \frac{\lambda}{2} \|\bm{w}\|^2_2+\sum_{i\in \bar{\mathcal{N}}_t}{\left(\max\left(1-y_i(\bm{w}^{\top}x^i-b),0\right)\right)^2}\right\}. \\
\end{split}
\label{model_2}
\end{align}


\subsubsection{The Inexact Bundle Method}

\textcolor{red}{ab hier: Theorie fehlt\\
!!! notation - oder in prelininaries einfugen}
%ingredients for bundle method

To solve the given bilevel problem with the above presented nonconvex inexact bundle algorithm the algorithm jumps between the two levels. Once the inner level problems are solved for a given \(\lambda\) - this is possible with any QP-solver as the problems are convex - the bundle algorithm takes the outcome \(w\) and \(b\) and optimizes the hyper-parameter again.

The difficulty with this approach is that the bundle algorithm needs one subgradient of the outer level objective function with respect to the parameter \(\lambda\). However to compute this subgradient also one subgradient of \(w\) and \(b\) with respect to \(\lambda\) has to be known.

\textcolor{red}{example in differentiable case}\\
Let us first assume that the outer and inner objective functions and \(w(\lambda)= \argmin \mathcal{L}_{low}(w,\lambda)\) are sufficiently often continuously differentiable to demonstrate the procedure of calculating the needed {(sub-)gradients}.

Let \(\mathcal{L}_{upp}(w,\lambda)\) be the objective function of the outer level problem, where the variable \(b\) was left out for the sake of simplicity.
To find an optimal hyper parameter \(\lambda\) given the input \(w\) the gradient \(g^{upp}_{\lambda}\) of \(\mathcal{L}_{upp}\) with respect to \(\lambda\) is needed in every iteration of the solving algorithm.
In order to calculate this gradient the chain rule is used yielding

\begin{equation*}
	g^{upp}_{\lambda} = \left(\frac{\partial}{\partial w}\mathcal{L}_{upp}(w,\lambda) \right)^{\top}\frac{\partial w(\lambda)}{\lambda}+\frac{\partial}{\partial \lambda}\mathcal{L}_{upp}(w,\lambda).
\end{equation*}

The challenge is here to find the term \(\frac{\partial w(\lambda)}{\lambda}\) because 

\[ \frac{\partial w}{\partial \lambda} \in \frac{\partial }{\partial \lambda} \argmin_{w} \mathcal{L}_{low}(w,\lambda). \]

Assuming \(\mathcal{L}_{low}\) is twice continuously differentiable in the optimality condition 
\begin{equation*}
	0 = \frac{\partial}{\partial w}\mathcal{L}_{low}(w^*,\lambda)
\label{opt_con}
\end{equation*}

opt point \(w^*\)

can be used to calculate the needed gradient in an indirect manner.

For these calculations to be possible the inner level loss function must yield a linear optimality condition in \(w\). This is for example the case for SVM loss functions with a squared one- or two-norm.
The optimality condition can then be written as the linear system

\begin{equation*}
	H(\lambda)w = h(\lambda).
\end{equation*}

By taking the partial derivative with respect to \(\lambda\) on both sides of the system one gets

\begin{equation*}
	\left( \frac{\partial H(\lambda)}{\partial \lambda}\right)w+H(\lambda)\frac{\partial w}{\partial \lambda} = \frac{\partial h(\lambda)}{\partial  \lambda}.
\end{equation*}

If \(H(\lambda)\) is invertible for all \(\lambda \in \Lambda\) then the needed gradient is given by 

\begin{equation*}
	\frac{\partial w}{\partial \lambda} = H^{-1}(\lambda)\left(\frac{\partial h(\lambda)}{\partial \lambda}-\left(\frac{\partial H(\lambda)}{\partial \lambda}\right)w\right).
\end{equation*}

\textcolor{red}{why H invertible??? - because we assume existence of \(w\)???}\\

\textcolor{red}{now for subgradients} \\

In practice we cannot expect \(\mathcal{L}_{low}\) to satisfy such strong differentiability properties.

\textcolor{red}{what kind of differentiability do I assume??? -> if I use normal Hinge then only locally Lipschitz $L_in$}

\begin{itemize}
	\item notation
	\item definition of subgradient-''matrix''
	\item chain rule
	\item optimality condition
\end{itemize}

\begin{itemize}
	\item welche art von inexaktheit -> Funktionswerte \(w,b\) inexakt\\
	-> gradient im Endeffekt exakt, da von exakter optimalit'tsbedingung ausgegangen wird
\end{itemize}

An important result about Lipschitz functions is Rademacher's theorem which states that these functions are differentiable almost everywhere but on a set of Lebesgue measure zero\cite[Theorem 3.1]{Heinonen2004}. 
Clarke deduces from this that the subdifferential at each of the nondifferentiable points is the convex hull of the limits of the sequence gradients a these points \cite[see Theorem 2.5.1]{Clarke1990}.

In practice this means that it is possible to choose a subgradient by using the (one sided) gradients at nondifferentiable points.
We keep this in mind when analyzing the procedure of finding a subgradient \( g^{\lambda} \in \partial^{\lambda} w(\lambda)\) in the nondifferentiable case.

\textcolor{red}{Notation}\\

To facilitate readability we use the following notation for the derivation of the nondifferentiable results.

The 'partial' subdifferential of a function \(f(a^*,\bar{b},\bar{c},...)\) at the point \(a^*\) with respect to one variable \(a\) when all other variables are fixed is denoted by 
\begin{equation*}
	\partial^{a} f(a^*,\bar{b},\bar{c},...).
\end{equation*}

A subgradient of this subdifferential is written \(g^{a} \in \partial^{a} f(a^*,\bar{b},\bar{c},...)\).

\textcolor{red}{what else???}

\textcolor{red}{explanation}





To calculate a subgradient \\
Chain rule for subdifferential \\
before: definition of subgradient in \(\R^m\)
\begin{theorem}[c.f. {\cite[Theorem 7.1]{Rockafellar1985}}]
	Let \(p(x) = f(F(x))\), where \(F:\R^n \to \R^d\) is locally Lipschitz and \(f:\R^d \to \R\) is lower semicontinuous. 
	Assume \[\nexists y \in \partial^{\infty}f(F(\bar{x})), y \neq 0 \quad \text{with} \quad 0 \in y \partial F(\bar{x}). \]
	Then for the sets 
	\[ M(\bar{x}):= \partial f(F(\bar{x}))\partial F(\bar{x}), \quad M^{\infty}(\bar{x}):= \partial^{\infty}f(F(\bar{x}))\partial F(\bar{x}), \]
	one has \(\hat{\partial}p(\bar{x}) \subset M(\bar{x})\) and \(\hat{\partial}^{\infty} p(\bar{x}) \subset M^{\infty}(\bar{x})\).
\end{theorem}



For me: \(f\) locally Lipschitz??? then partial derivatives are the same! Else: check definition of derivatives!

\textcolor{red}{-> theory partial derivatives for subgradients????????? \\
??? Formula \(??? \in \partial L_{upp}{\partial \lambda}\) \\
???one has to assume that the inner level problem is locally Lipschitz (or more general: its nonconvex subdifferential is well defined at every point). \\
Subdifferential has to have again a subdifferential!!! -> w.r.t. \(\lambda\)}


The main idea is to replace the inner level problem by its optimality condition


\(\partial(w,b)\) means in this case that the subdifferential is taken with respect to the variables \(w\) and \(b\). \\
-> theory for subdifferentials in more than one variable!!! \\

For convex inner level problem this replacement is equivalent to the original problem.

The difference to the approach described in \cite{Kunapuli2008} is that the problem is not smoothly replaced by its KKT conditions but only by this optimality condition. The weight vector \(\bm{w}\) and bias \(b\) are treated as a function of \(\lambda\) and are optimized separately from this hyper-parameter.
The reformulated bilevel problem becomes:

\begin{align}
	\begin{split}
	\min_{\bm{W},\bm{b}} \quad &  \mathcal{L}_{hinge}(\bm{W},\bm{b}) = \frac{1}{T}\sum_{t=1}^T\frac{1}{|\mathcal{N}_t|}\sum_{i \in \mathcal{N}_t}{\max\left(1-y_i((\bm{w}^t)^{\top}x-b_t),0\right)}\\
	\text{subject to} \quad & \lambda > 0 \\
	& \text{for } t = 1,...,T \\
	& 0 \in \partial(w,b)\mathcal{L}_{low}(\lambda,w^t,b_t) \\
\end{split}
\label{SVM_opt_cond}
\end{align}

where \(\mathcal{L}_{low}\) can be the objective function of either of the two presented lower level problems.



solve the inner level problem (quadratic problem in constrained case) by some QP solver \\
put solution into upper level problem and solve it by using bundle method \\
difficulty: subgradient is needed to build model of the objective function --> need subgradient \(\frac{\partial \mathcal{L}}{\partial \lambda}\) --> for this need \(\frac{\partial (W,b)}{\partial \lambda}\) \\
but \((w,b)\) not available as functions -> only values

Moore et al. \cite{Moore2011} describe a method for getting the subgradient from the KKt-conditions of the lower level problem:

lower level problem convex -> therefore optimality conditions (some nonsmooth version -> source???) necessary and sufficient -> make ``subgradient'' of optimality conditions and then derive subgradient of w, b from this. \\
---> what are the conditions? optimality condition Lipschitz? 

Say (show) that all needed components are locally Lipschitz; state theorems about differentiability almost everywhere and convex hull of gradients gives set of subgradients\\
introduce special notation (only for this chapter) and because of readability adopt ``gradient writing''

Subgradients:
\(\mathcal{G}_{upp,\lambda}, \mathcal{G}_{upp,w},\mathcal{G}_{upp,b}\) -> subgradients of outer objective \\
\(g_w, g_b\) -> subgradient of w, b

\[ final subgradient = \left(\mathcal{G}_{upp,w}(w,b,\lambda)\right)^{\top}g_w+\left(\mathcal{G}_{upp,b}(w,b,\lambda)\right)^{\top}g_b+\mathcal{G}_{upp,\lambda}(w,b,\lambda) 
\label{subgr_upp}\]

subgradients \(\mathcal{G}_{upp,...}\) easy to find (assumption that locally Lipschitz) -> in this application differentiable

difficulty: find \(g_w, g_b\)
important: optimality condition must be a linear system in \(w,b\) -> this is the case in this application
\[ H(\lambda)\cdot (w,b)^{\top} = h(\lambda) \]

find subgradients of each element (from differentiation rules follows)

\[ \partial_{\lambda} H\cdot (w,b)^{\top} + H \cdot(\partial_{\lambda} w, \partial_{\lambda} b)^{\top} = \partial_{\lambda} h  \]

solve this for \((w,b)\):
\[ (\partial_{\lambda} w, \partial_{\lambda}b)^{\top} = H^{-1}\left(\partial_{\lambda}h-\partial_{\lambda} H \cdot(w,b)^{\top}\right) \]

matrix \(H\) has to be inverted -> in the feature space so scalable with size of data set -> still can be very costly \cite{Moore2011}

Applied to the two bilevel classification problems derived above, the subgradients have the following form:

derivative of upper level objective:
Notation: \(\delta_i := 1-y_i(w^{\top}x^i-b)\)

\begin{align}
	\partial_{w}\mathcal{L}_{upp}&= \frac{1}{T}\sum_{t=1}^T\frac{1}{\mathcal{N}_t}\sum_{i \in \mathcal{N}_t}{\left\{\begin{array}{cl} -y_ix^i & \text{if } \delta_i >0 \\ 0 & \text{if } \delta_i \leq 0 \end{array} \right.} \\
	\partial_{b}\mathcal{L}_{upp}&= \frac{1}{T}\sum_{t=1}^T\frac{1}{\mathcal{N}_t}\sum_{i \in \mathcal{N}_t}{\left\{\begin{array}{cl} y_i & \text{if } \delta_i >0 \\ 0 & \text{if } \delta_i \leq 0 \end{array} \right.}
	\end{align}
	
	here at the kink subgradient 0 is taken

for hingequad: -> here subgradient \\
optimality condition:
\begin{align}
	0 = \partial_{\bm{w}}\mathcal{L}_{low} &= \lambda \bm{w}+2\sum_{i \in \bar{\mathcal{N}_t}}{\left\{\begin{array}{cl} (1-y_i(w^{\top}x^i-b))(-y_ix^i) & \text{if } \delta_i >0 \\ 0 & \text{if } \delta_i \leq 0 \end{array} \right.} \\
	0 = \partial_b\mathcal{L}_{low} &= 2\sum_{i \in \bar{\mathcal{N}_t}}{\left\{\begin{array}{cl} (1-y_i(w^{\top}x^i-b))(y_i) & \text{if } \delta_i >0 \\ 0 & \text{if } \delta_i \leq 0 \end{array} \right.}
	\end{align}
	
subgradient??? is this smooth? with respect to \(\lambda\)
\begin{align}
	0 &= \bm{w}+\lambda\partial_{\lambda}\bm{w}+2\sum_{i \in \bar{\mathcal{N}_t}}{\left\{\begin{array}{cl} (-y_i(\partial_{\lambda}w^{\top}x^i-\partial_{\lambda}b))(-y_ix^i) & \text{if } \delta_i >0 \\ 0 & \text{if } \delta_i \leq 0 \end{array} \right.} \\
	0 &= 2\sum_{i \in \bar{\mathcal{N}_t}}{\left\{\begin{array}{cl} (-y_i(\partial_{\lambda}w^{\top}x^i-\partial_{\lambda}b))(y_i) & \text{if } \delta_i >0 \\ 0 & \text{if } \delta_i \leq 0 \end{array} \right.}
\end{align}

From this the needed subgradients can be calculated via:

\begin{equation}
	2\cdot\begin{pmatrix} \sum_{i \in \bar{\mathcal{N}_t}}\frac{\lambda}{2}+{y_i^2x^i(x^i)^{\top}} & \sum_{i \in \bar{\mathcal{N}_t}}	{-y_i^2x^i} \\ \sum_{i \in \bar{\mathcal{N}_t}}{-y_i^2(x^i)^{\top}} & \sum_{i \in \bar{\mathcal{N}_t}}{y_i^2}\end{pmatrix}
	\cdot \begin{pmatrix} \partial_{\lambda}w \\ \partial_{\lambda}b\end{pmatrix}
	= \begin{pmatrix}-w \\ 0\end{pmatrix}
\label{subgr_wb}
\end{equation}

for hinge not quad:

not as much information in the subgradient/derivative

similar calculation leads to 

\begin{align}
	\partial_{\lambda}w &= -\frac{w}{\lambda} \\
	\partial_{\lambda}b &= 0
\end{align}

\subsubsection{The Algorithm???}

The inexact bundle algorithm for the support vector classification task in bilevel formulation 
\vspace{1.5em}
\hrule  \vspace{0.4ex} \hrule
\vspace{1ex}
\textbf{Bilevel Bundle Method}
\vspace{1ex}
\hrule
\vspace{1ex}
Initiate all parameters, select a starting hyper-parameter \(\lambda_1\) and solve the lower level problem for \(\bm{w}^1\) and \(b_1\). \\
Calculate arbitrary subgradients of \(\bm{w}^1\) and \(b_1\) with respect to \(\lambda\) via \ref{subgr_wb} and a subgradient of the upper level problem by \ref{subgr_upp}.  
For \(k = 1,2,3,  \dotsc \)   

\begin{enumerate}
	\item Calculate the step \(d^k\) by minimizing the model of the convexfied objective
	\item Compute the aggregate subgradient and error and the stopping tolerance \(\delta\).
		If \(\delta_k \leq \mathtt{tol} \rightarrow \) STOP.
	\item Set \( \lambda^{k+1} = \hat{\lambda}^k + d^k \).
	\item solve again the inner level problem and calculate all subgradients needed to compute a subgradient of the outer level objective \\
	Calculate function value and a subgradient for the outer level objective function and test if a serious step was done 
	If yes, set \(\hat{\lambda}^{k+1} = \lambda^{k+10}\) and select \(t_{k+1} > 0\). \\
	Otherwise \(\rightarrow\) nullstep \\
	Set \(\hat{\lambda}^{k+1} = \hat{\lambda}^k\) and choose \(0 < t_{k+1} \leq t_k\). 	
	\item Select new bundle index set \(J_{k+1}\). Calculate convexification parameter \(\eta_k\)
	and update the model \(M^k\)
\end{enumerate}
\vspace{1ex}
\hrule
\vspace{1.5em}

Names for algorithms: BBMH -> hinge as inner level, BBMH2 -> hingequad as inner level 

\subsection{Numerical Experiments}

The bilevel-bundle algorithm for classification was tested for four different data sets taken from the UCI Machine Learning Repository \emph{citations as said in ``names'' data??? }.
For comparability with the already existing results presented in \cite{Kunapuli2008} the following data and specifications of it were taken:

\begin{center}
\emph{Table like in Kunapuli}
\begin{table}[H]%
	\begin{tabular}{lcccc}
		\hline
    Data set & \(l_{train}\) & \(l_{test}\) & n & T \\
		\hline
		Pima Indians Diabetes Database & 240 & 528 & 8 & 3 \\
		Wisconsin Breast Cancer Database & 240 & 443 & 9 & 3 \\
		Cleveland Heart Disease Database & 216 & 81 & 13 & 3 \\
		John Hopkins University Ionosphere Database & 240 & 111 & 33 & 3
	\end{tabular}
	\caption{}
	\label{}
\end{table}
\end{center}

As described in the PhD thesis the data was first standardized to unit mean and zero variance (\emph{not the 0,1 column in ? dataset}). The bilevel problem with cross validation was executed 20 times to get averaged results.
The results are compared by cross validation error, test error -> write which error this is and computation time.
Additionally write \(\bm{w}\), \(b\), \(\lambda\) ???
The objective function and test error were scaled by 100. -> also test error (to get percentage) 

After every run the calculated \(\lambda\) was taken and the algorithm was trained with \(\frac{T}{T-1} \lambda\) on the whole training set.
Then the percentage of misclassifications on the test set was calculated via

\begin{equation}
	E_{test} = \frac{1}{l_{test}}\sum_{i=1}^{l_{test}}{\frac{1}{2}|sign\left(\bm{w}^{\top}x^i-b \right)-y_i|}
\label{test_err}
\end{equation}

Table ??? shows the results 

\begin{center}
\begin{table}[H]%
	\begin{tabular}{lllll}
		\hline
		Data set & Method & CV Error & Test Error & Time (sec.) \\
		\hline
		\texttt{pima} & hingequad & \(60.72 \pm 9.56\) & \(24.11\pm 2.71\) & \(2.15 \pm 0.52\)\\
		 &              hinge loss & & & \\
		\texttt{cancer} &  hingequad & \(10.75\pm 7.52\) & \(3.41 \pm 1.16\) & \(3.43 \pm 28.84\) \\
		 &              hinge loss & & & \\
		\texttt{heart} &  hingequad & \(48.73 \pm 5.53\) & \(15.56 \pm 4.44\) & \(3.43 \pm 43.39\)\\
		 &              hinge loss & & & \\
		\texttt{ionosphere} &  hingequad & \(39.30 \pm 5.32\) & \(12.21 \pm 4.10\) & \(14.17 \pm 51.27\)\\
		 &              hinge loss & & & \\
	\end{tabular}
	\caption{}
	\label{res_table}
\end{table}
\end{center}

%\textcolor{red}{interesting: 0 computing time for ionosphere??? - solved now -> vertippt}

Extra table for \(\bm{w}\), \(b\), \(\lambda\) ?

First experiment: Classification

Write down bilevel classification problem and (if needed) which specification of the inexact bundle algorithm is used.

%Write down the sets were used and how they were prepared.















\textcolor{red}{\textbf{Covtype tests}\\
Datensatz zerteilt: 1000, 5000, 10000, 50000 Datens\"atze\\
Ergebnisse mit Matlab App (linear SVM, 3 folds, parallel used):\\
\begin{tabular}{c|c|c}
Datensatz & Zeit & Fehler \\
\hline
1,000 & 16.22 sek & 34.2\% \\
5,000 & 10.524 sek & 18.3\% \\
10,000 & 14.689 sek & 16.5\% \\
50,000 &  643.57 sek & 16.9\% \\
50,000 (Rechnerhalle, 4 parallel) & 326.83 & 16.9\% \\
100,000 --''-- & 2492 sek = 41.5333 min & 21.3\%
\end{tabular}}

\textcolor{red}{scheint bei 50000 tats\"achlich so lange zu dauern, da ziemlich genau doppelt so schnell bei parallel-Rechnung mit 4 anstatt 2 ``Rechnern'' --> kein Arbeitsspeicher Problem}

\textcolor{red}{Test mit bundle bilevel-Algorithmus:\\
for covtype, ``Hare''-stopping condition\\
\begin{tabular}{c|c|c|c|c|c|c}
trainigs set & starting value & lambda & time & k & inull & Fehler\\
\hline
1,000 & 1 & 86.3493 & 62.8280 sek & 29 & 0 & \\
1,000 & 86 & 86.3488 & 130.1105 sek & 61 & 8 & 19\% \\
1,000 (quadprog) & 86 & 86.3488 & 6.5572 sek & 61 & 8 & \\
5,000 qp & 1 & 46.4142 & 16.2338 sek & 27 & 0 & \\
5,000 qp & 47 & 46.4142 & 34.9573 sek & 60 & 0 & 14.74\% \\
10,000 qp & 1 & 49.8146 & 59.7462 sek & 49 & 0 & \\
10,000 qp & 50 & 49.8146 & 85.0816 sek & 69 & 0 & 15.32\% \\
50,000 qp & 58.6842 & 588.2828 sek & 45 & 0 & \\
50,000 qp & 60 & 58.6842 & 897.8123 sek & 69 & 0 & 14.86\% \\
100,000 qp & 1 & 73.3711 & 1358.7 sek = 22.6455 min & 37 & 0 & 39.41
\end{tabular}
\\
Fehler berechnet f\"ur Daten 1001 bis 2000 von komplettem Datensatz\\
Versuch mit Daten 1001-5000: scheint Problem bei matrizenaufbau zu haben \\
Analyse Timer: 100\% der Zeit in postpro_wb_class_hinge_qpas dabei 100\% der Zeit für qpas\\
Ergebnis viel!!! schneller, wenn quadprog und sparse-matrizen, gibt EXAKT! gleiches Ergebnis für fehler\\
geht dann auch mit mehr datensätzen (4000) -> Fehler nur 7.75\%??? \\
f\"ur Datensatz 5000 \(\to\) Testdaten 5001 bis 10000 \\
exemplarisch getestet: wie groß ist der Unterschied bei ergebnisse bei verschiedenen Starwerten? - bei 10000: 1e-15 , bei 50000: exact\\
f\"ur Datensatz 10000: Testdaten 10001 bis 15000 \\
50000: deutlicher Anstieg der Rechenzeit merkbar irgendwo zwischen 10000 und 50000 muss eine Schwelle liegen - Speicher? - ab dieser Matrix Gr\"osse gibt matlab fehler wenn nicht sparse matrizen explizit erstellt werden sollen (zu viel Speicher) unwahrscheinlich - siehe erkl\"arung App \\
komish: warum dauert es bei näherem Startwert so viel länger??? }





\addtocontents{toc}{%
  \protect\contentsline{section}{References}{}}
\singlespacing
\bibliography{Bibliography}
	\bibliographystyle{plain}


\end{document}